\documentclass{ieeeaccess}
\usepackage{cite}
\usepackage{amsmath,amssymb,amsfonts}
\usepackage{algorithmic}
\usepackage{graphicx}
\usepackage{textcomp}
\def\BibTeX{{\rm B\kern-.05em{\sc i\kern-.025em b}\kern-.08em
    T\kern-.1667em\lower.7ex\hbox{E}\kern-.125emX}}
\begin{document}
\history{Date of publication xxxx 00, 0000, date of current version xxxx 00, 0000.}
\doi{10.1109/TQE.2025.DOI}
 
\title{Text Encoding in Quantum Natural Language Processing: A Systematic Survey}

\author{\uppercase{Danilo Cavaliere}\authorrefmark{1}, \uppercase{Giuseppe Fenza}\authorrefmark{1}, \uppercase{Lukasz Gajewski}\authorrefmark{1}, and \uppercase{Vincenzo Loia}\authorrefmark{1}}

\address[1]{Department of Management \& Innovation Systems, Università degli Studi di Salerno, Italy}

\markboth{TEXT ENCODING IN QUANTUM NATURAL LANGUAGE PROCESSING}{TEXT ENCODING IN QUANTUM NATURAL LANGUAGE PROCESSING}

\corresp{Corresponding authors: Danilo Cavaliere and Lukasz Gajewski (e-mails: \{dcavaliere, lgajewski\}@unisa.it).}

\begin{abstract}
Quantum natural language processing (QNLP) applies superposition, entanglement, and interference to the representation and processing of linguistic information. Any QNLP pipeline, whatever its architecture, must first map discrete or distributed text representations onto qubit registers or parameterized-circuit inputs. This encoding step governs model expressivity, qubit resource requirements, and trainability. Yet it is rarely studied as an independent variable, and no unified method exists for analyzing and comparing encodings. This article closes that gap with a systematic literature review of text encoding strategies in QNLP. The review covers 125 primary studies published between 2021 and 2026, under a protocol defined before the search was executed that fixed the search strategy, eligibility criteria, screening procedure, snowball sampling, and data extraction. It introduces a unifying taxonomy for analyzing and comparing encodings. Seven encoding families are identified and characterized: four major families, angle, amplitude, syntactic (DisCoCat), and structured feature maps (IQP/ZZ), and three emerging families: data re-uploading, hybrid classical--quantum embedding, and Hamiltonian/observable encoding. Their distribution is analyzed across NLP tasks, simulation and hardware environments, qubit budgets, and reporting practice. Three findings stand out: noiseless circuit simulation dominates the literature, real-hardware execution remains rare, and evaluation is too fragmented across author-chosen datasets, baselines, and metrics for results to be compared between studies. The encoding bottleneck, the dimensionality mismatch between linguistic representations and near-term qubit capacities, is stated explicitly as a basis for resource-efficient QNLP algorithm and hardware design.
\end{abstract}

\begin{keywords}
Encoding bottleneck, parameterized quantum circuits, quantum machine learning, quantum natural language processing, systematic literature review, text encoding.
\end{keywords}

\titlepgskip=-15pt

\maketitle

\section{Introduction}
\label{sec:introduction}
\PARstart{E}{very} quantum natural language processing (QNLP) system begins with a mandatory step: converting text into a quantum state. Before a single gate is applied or a single parameter is trained, a token, sentence, or document must be mapped onto qubits. This mapping, the encoding, determines how many qubits are consumed and how deep the resulting circuit must be. It determines whether gradients can be computed efficiently, and whether any genuinely quantum property, such as entanglement, is brought to bear on the linguistic content. It is, in short, the most consequential design decision in the pipeline.

Despite this centrality, encoding is rarely the subject of dedicated study. Most QNLP papers dispatch it in a paragraph, introduce one strategy without justification, and move on to reporting accuracy numbers. After five years of accelerating publication activity, the result is a fragmented body of work: dozens of encoding families coexist, yet there is no common approach for analyzing and comparing them. Unified taxonomies, systematic comparisons, and a shared understanding of encoding trade-offs are all missing from the literature. Existing QNLP surveys and reviews do not isolate encoding as a variable \cite{chen_survey_2024,pallavi_quantum_2025}. The few explicitly comparative studies contrast model families or ans\"atze rather than encoding schemes \cite{bouakba_comparative_2025,castillo_comparative_2025}. This article fills that gap: it conducts an extensive systematic literature review of the topic and makes text encoding in QNLP the primary subject of study.

\subsection{Scope}
This survey reviews primary research that (1) applies a genuine quantum circuit (simulated noiseless, simulated noisy, or executed on real hardware) to at least one NLP task, and (2) identifies, at least by name or type, the encoding scheme used to map text or text-derived features onto qubits or circuit parameters. The review follows an SLR protocol defined before the search (Section~\ref{sec:method}) adapted from established guidelines for secondary studies in computer science \cite{keele2007guidelines}.

Papers that use quantum-inspired or quantum-motivated classical algorithms without instantiating an actual parameterized quantum circuit are excluded. Also excluded are survey and review papers, hardware-focused benchmarking studies that treat NLP datasets as a generic payload, and papers with insufficient methodological detail to assess the encoding scheme. The final corpus comprises 125 primary studies.

The scope is limited to the encoding step. Quantum circuit design, ansatz expressibility, and quantum error correction are surveyed only where they interact directly with encoding choices. Classical NLP methods are discussed only as baselines.

\subsection{Key Contributions}
This article makes four contributions to the study of text encoding in QNLP:
\begin{enumerate}
\item[1)] A novel taxonomy of text encoding strategies, organized by the point at which text enters the quantum computation and covering seven encoding families (four major and three emerging) with their specialized variants.
\item[2)] A systematic review of 125 primary studies that apply or propose quantum text encodings across NLP tasks, conducted under a protocol defined before the search.
\item[3)] A quantitative cross-study characterization of the corpus using the classification scheme defined in Section~\ref{subsec:analysis_framework}, spanning encoding-by-model-family, temporal, task-fit, and experimental-rigor analyses.
\item[4)] A corpus-derived set of five open challenges, including an explicit statement of the encoding bottleneck and an actionable agenda for future work.
\end{enumerate}

\subsection{Related Work}
\label{subsec:related_work}
Several prior works survey the QNLP literature, but none treats text encoding as the central variable. Varmantchaonala et al.\ \cite{varmantchaonala_quantum_2024} provide the most comprehensive QNLP survey to date, covering quantum algorithms, circuit architectures, and NLP task categories; encoding appears only as an implementation detail. Guarasci et al.\ \cite{guarasci_quantum_2022} offer an early and influential overview of QNLP challenges and opportunities, from theoretical quantum-advantage arguments to hardware experiments, without isolating encoding as a dimension of comparison. Widdows et al.\ \cite{widdows_quantum_2024} examine word embeddings and sequential quantum models with more attention to representation trade-offs than any other general survey, but stop short of a systematic comparison. Bouakba and Belhadef \cite{bouakba_comparative_2025} experimentally compare compositional DisCoCat models against quantum neural networks; their compositional-versus-flat-vector axis overlaps with the taxonomy proposed here, though their scope is limited to two model families. Huzzat et al.\ \cite{huzzat_primer_2026} is the most encoding-aware prior work: it addresses the mismatch between text-representation dimensionality and the qubit capacity of NISQ classifiers, and treats classical compression as a prerequisite for stable training. It remains, however, a conceptual primer for a single model class rather than a systematic review.

Beyond QNLP itself, the quantum engineering literature has an established tradition of comprehensive subfield surveys, covering, for instance, quantum circuit compilation for NISQ devices \cite{kusyk2021survey}, quantum computing for finance \cite{egger2020quantum}, and entanglement routing in quantum networks \cite{abane2025entanglement}, the last of which likewise organizes its field through an explicit taxonomy of approaches. None of these addresses text encoding in QNLP, which is the gap the present survey fills within that tradition.

This survey complements these works by treating encoding as the primary unit of analysis, yielding the first quantitative cross-study characterization of how text is encoded in QNLP and at what resource cost.

\subsection{Structure}
The rest of the article is organized as follows. Section~\ref{sec:background} provides the background on quantum computing, compositional distributional semantics, and the role of encoding in QNLP. Section~\ref{sec:method} describes the systematic review protocol (search strategy, eligibility criteria, screening, and snowball sampling) and the data extraction and classification scheme used throughout the analysis. Section~\ref{sec:taxonomy} introduces the taxonomy of text encoding strategies and characterizes each family. Section~\ref{sec:comparative_analysis} analyzes how each encoding family is used across model families and ans\"atze. Section~\ref{sec:temporal_evolution} traces the temporal evolution of encoding and model choices. Section~\ref{sec:task_fit} examines the fit between encoding, model family, and NLP task. Section~\ref{sec:experimental_rigor} assesses experimental rigor across simulation regimes, hardware use, and reproducibility. Section~\ref{sec:trends_challenges} synthesizes the trends, states the encoding bottleneck explicitly, and enumerates the open challenges and the future-work agenda. Section~\ref{sec:conclusion} concludes.

\section{Background}
\label{sec:background}

\subsection{Quantum Computing Primitives}

A quantum bit (qubit) occupies a superposition of the basis states $|0\rangle$ and $|1\rangle$:
\begin{equation}
|\psi\rangle = \alpha|0\rangle + \beta|1\rangle, \quad |\alpha|^2 + |\beta|^2 = 1,
\end{equation}
where $\alpha, \beta \in \mathbb{C}$ are probability amplitudes. A register of $n$ qubits occupies a $2^n$-dimensional Hilbert space $\mathcal{H} = \mathbb{C}^{2^n}$; its state is a unit vector in that space. This exponential dimension is the property every amplitude-style encoding exploits. Entanglement introduces correlations with no classical counterpart, and measurement collapses the state to a classical bitstring, destroying superposition. Whether superposition, entanglement, or interference confers a practical benefit on NLP tasks remains the field's central open question, and Section~\ref{sec:experimental_rigor} examines how far the surveyed evidence can speak to it.

\subsection{Parameterized Quantum Circuits}

The dominant computational model in near-term quantum machine learning is the parameterized quantum circuit (PQC), also called a variational quantum circuit \cite{cerezo2021variational}. A PQC applies a sequence of unitary gates to an initial state $|\psi_0\rangle$:
\begin{equation}
U(\boldsymbol{\theta}) = U_L(\boldsymbol{\theta}_L) \cdots U_1(\boldsymbol{\theta}_1),
\end{equation}
where $\boldsymbol{\theta} \in \mathbb{R}^p$ are trainable parameters. Model outputs are expectation values $\langle O \rangle = \langle \psi_0 | U^\dagger(\boldsymbol{\theta})\, O\, U(\boldsymbol{\theta}) | \psi_0 \rangle$, and gradients follow analytically from the parameter-shift rule \cite{mitarai2018quantum, schuld2019evaluating}. PQCs underlie six of the seven encoding families surveyed here; the exception is grammatical encoding, which derives circuit structure from sentence syntax rather than from a fixed ansatz.

PQCs operate under the constraints of noisy intermediate-scale quantum (NISQ) hardware \cite{preskill2018quantum}: shallow circuits, limited qubit counts, and significant gate error rates. These constraints bear directly on the encoding step. Every qubit or gate layer spent loading a text representation is a resource unavailable to the learning task.

\subsection{Text Encoding}

Text encoding is the procedure that maps a text input $x$ (a token, sentence, or document) to a quantum state $|\phi(x)\rangle \in \mathcal{H}$ or to a parameter vector $\boldsymbol{\theta}(x)$ that initializes a PQC. It is the mandatory interface between classical language data and the quantum circuit. The input $x$ may be a raw token sequence, a classical embedding vector from a pre-trained model, or a parse tree carrying grammatical structure. Which representation is used, and how it is loaded onto qubits, defines the encoding family.

\subsection{Compositional Distributional Semantics and DisCoCat}

One of the two main traditions in QNLP originates in categorical linguistics. Coecke, Sadrzadeh, and Clark \cite{coecke2010mathematical} showed that the compact closed categorical structure of pregroup grammars also governs quantum processes in Hilbert space. In a pregroup grammar, each word carries a grammatical type, and a sentence is parsed by reducing compound types through tensor contractions. These reductions obey the same diagrammatic calculus as quantum processes, so a parse tree translates directly into a quantum circuit: words become parameterized quantum states, and grammatical reductions become cups in a string diagram.

This is the DisCoCat (Distributional Compositional Categorical) model, operationalized as an end-to-end pipeline by \cite{lorenz_qnlp_2023} and executed on near-term hardware in \cite{meichanetzidis_quantum_2021, meichanetzidis_grammar-aware_2023} via the \textsc{lambeq} toolkit. The circuit structure is not fixed in advance: the parse tree determines it, so syntactically distinct sentences produce structurally distinct circuits. This property sets the grammatical encoding family apart from all other families, which treat the sentence as a flat feature vector.

\subsection{Quantum Machine Learning for Text}

The second tradition approaches NLP as a classification problem. It applies quantum feature maps \cite{schuld2019quantum} and variational circuits \cite{cerezo2021variational} to text-derived feature vectors: classical embeddings (word2vec, GloVe, or transformer outputs) are compressed to the available qubit count, loaded into a PQC via rotation gates or amplitude encoding, and trained end-to-end on labeled examples. Because this approach assumes nothing about grammatical structure, it applies to arbitrary text without a parser and scales more easily to realistic dataset sizes. Its logical endpoint is the hybrid quantum--classical architecture, in which a large classical language model handles high-dimensional representation and a small quantum circuit acts as a downstream head; such architectures have grown in prevalence since 2023.

The two traditions differ not only in how they encode text but in what they claim the quantum component contributes. Compositional models argue for structural alignment between grammar and quantum computation; QML models argue for expressivity or kernel separability.

\section{Review Methodology}
\label{sec:method}

\Figure[t!](topskip=0pt, botskip=0pt, midskip=0pt)[width=\textwidth]{PRISMA_Chart.jpg}{PRISMA flow diagram. Final corpus: $n = 125$ (95 from database search, 30 from snowball sampling).\label{fig:prisma}}

This systematic literature review followed the guidelines for secondary studies in software engineering established by Keele and Charters \cite{keele2007guidelines}. The protocol was defined before the search was executed. It fixed the search strategy, eligibility criteria, screening procedure, snowball sampling, and data extraction process.

\subsection{Search Strategy}
Two complementary search queries were executed on two electronic databases, Scopus and IEEE Xplore, to cover both the quantum machine learning and the compositional QNLP literature. The first query targeted the intersection of quantum computing or machine learning with NLP and encoding-related terms. The second targeted QNLP-specific terminology: the DisCoCat and DisCoCirc frameworks, established pipelines that implement them such as the \textsc{lambeq} toolkit, and categorical linguistic concepts. The searches were restricted to English-language journal articles and conference papers published between 2021 and 2026 (ends included), excluding editorials, book chapters, and erratum notices. They returned 269 results from Scopus and 217 from IEEE Xplore, a combined total of 486 records. Removing 110 duplicates left 376 unique records.

\subsection{Eligibility Criteria}
\label{sec:criteria}

A study was selected if and only if it met all six inclusion criteria below.

\begin{itemize}
    \item \textit{IC1: NLP task grounding.} The study applies or evaluates the proposed method on at least one natural language processing task, such as classification, similarity, sequence labeling, question answering, or generation.

    \item \textit{IC2: Genuine quantum computation.} The study implements and evaluates the method in at least a noiseless quantum circuit simulation (exact statevector simulation or equivalent) or on real quantum hardware; merely quantum-inspired classical algorithms, such as tensor-network models or classical approximations of quantum kernels, do not satisfy this criterion.

    \item \textit{IC3: Encoding scheme identified.} The study names or identifies at least the type of encoding used to map text or text-derived features onto qubits, quantum states, or circuit parameters; a named or typed mention is sufficient, and circuit diagrams or full mathematical derivations are not required.

    \item \textit{IC4: Peer-reviewed or preprint venue.} The study appeared in a peer-reviewed journal, conference, or workshop, or as a preprint with sufficient methodological detail to assess IC1--IC3.

    \item \textit{IC5: Language and access.} The study is available in full text, in English, and accessible for review.

    \item \textit{IC6: Time window.} The study was published between 2021 and 2026 inclusive.
\end{itemize}

Even if a study satisfied all six inclusion criteria, it was excluded if any of the following applied:
\begin{itemize}
    \item \textit{EC1: No NLP application.} The study applies quantum machine learning exclusively to non-text domains and offers no transferable text-encoding discussion.

    \item \textit{EC2: Classical-only encoding.} The study uses encoding purely in the classical sense (tokenization, positional encoding, one-hot) with no quantum circuit involved.

    \item \textit{EC3: Insufficient detail.} The study is an abstract, extended abstract, or short paper that names an encoding but provides no formal description or reproducible detail.

    \item \textit{EC4: Duplicate or superseded.} The study is a conference version superseded by an extended journal version; only the more complete version was retained.

    \item \textit{EC5: Pure hardware focus.} The study addresses hardware benchmarking alone and uses an NLP dataset only as a generic payload, with no discussion of linguistic encoding properties.

    \item \textit{EC6: Quantum-inspired, classical surrogate, or unverifiable simulation.} The study uses quantum-derived mathematical formalism without instantiating an actual parameterized quantum circuit, reports only classical surrogates of quantum behavior, or claims quantum simulation without evidence of framework or circuit use.

    \item \textit{EC7: Survey or review paper.} The study is primarily a secondary study rather than an original research contribution; such papers were retained as reference material and are discussed in Section~\ref{subsec:related_work}.
\end{itemize}

\subsection{Screening and Selection}

Screening proceeded in two stages. In the first stage, all 376 records were assessed on title and abstract alone for broad topical relevance to quantum computing applied to text. No IC/EC criteria were applied at this stage; its only purpose was to remove records clearly outside the scope of the review. This eliminated 234 records and left 142 reports sought for retrieval, of which 6 could not be retrieved, leaving 136 reports assessed for eligibility.

In the second stage, the full text of each of the 136 reports was assessed against the complete set of inclusion and exclusion criteria of Section~\ref{sec:criteria}. This yielded 95 papers for the primary corpus. The remaining 41 retrieved reports were excluded or retained as surveys, mainly under IC2 (no genuine quantum circuit; $n = 15$), IC3 (encoding scheme not identified; $n = 6$), EC6 (quantum-inspired, classical surrogate, or unverifiable simulation; $n = 3$), and EC7 (survey or review; $n = 11$); six records were excluded under other criteria.

\subsection{Snowball Sampling}
To reduce the risk of missing relevant work not indexed by the two databases, two rounds of backward and forward snowball sampling \cite{wohlin2014guidelines} were applied to the surveys excluded under EC7. Surveys were chosen as the seed set, rather than the included primary studies, for two reasons. First, Wohlin's guidelines ask for a start set that is small, diverse, and well connected to the target population; in a field as young as QNLP, a survey's reference list is by construction a curated, near-exhaustive sample of the primary literature up to its search date, so each survey seed contributes far more unique candidates than a primary study, whose citations concentrate on the same few foundational works. Second, forward snowballing from surveys exploits a stable citation practice: new primary studies routinely cite the field's reviews, so the papers citing a survey are enriched for precisely the recent primary work that database indexing lags behind. The cost of this choice is that a primary study neither cited by nor citing any survey of the field remains invisible to the procedure; this residual risk is acknowledged in Section~\ref{sec:conclusion}. Reference lists and citing papers yielded 1222 unique candidate records. Title-and-abstract screening for broad topical relevance reduced this set to 242 reports sought for retrieval, of which 37 could not be retrieved, leaving 205 reports assessed for eligibility. Full-text assessment against the complete IC/EC protocol then added 30 papers to the corpus. The remaining 175 retrieved reports were excluded or retained as surveys. The main grounds were IC1 (no NLP task; $n = 74$), IC2 (no genuine quantum circuit; $n = 52$), IC6 (outside the time window; $n = 20$), IC3 (encoding scheme not identified; $n = 4$), and EC7 (survey or review; $n = 22$); three records were excluded under other criteria.

\subsection{Final Corpus}
The final primary corpus comprises 125 papers, 95 from the database search and 30 from snowball sampling; Fig.~\ref{fig:prisma} summarizes the selection process as a PRISMA flow diagram.

\subsection{Data Extraction and Classification}
\label{subsec:analysis_framework}

Each primary study was classified along a fixed set of dimensions, so that every classification, cross-tabulation, and trend reported in Sections~\ref{sec:taxonomy}--\ref{sec:trends_challenges} follows directly from the extraction sheet rather than from an ad hoc reading of the corpus. The two primary dimensions are the \emph{encoding family}, defined by the point at which text or text-derived features enter the quantum computation (Section~\ref{sec:taxonomy}), and the \emph{model family}, defined by the architecture that consumes the encoded state. Secondary dimensions capture the resources and evidence attached to that choice: qubit count and whether qubit scaling was tested, execution regime, presence of a classical baseline, and whether scalability was discussed. Each dimension is operationalized as one or more columns of the structured extraction sheet compiled during data extraction.

Assignment follows explicit decision rules. A study is classified by its \emph{dominant} encoding interface: the one that carries the text representation into the circuit whose output is evaluated. Auxiliary encodings used only in ablations or baselines do not change the class; a study whose main pipeline is amplitude-encoded but that reports an angle-encoded ablation is classified as amplitude. Where a pipeline stacks a classical encoder in front of a quantum layer, the study is classified as \texttt{hybrid\_embed} only when the classical stage is co-trained or purpose-built as part of the contribution; a frozen pretrained embedding feeding a standard rotation layer remains \texttt{angle}. The \texttt{hybrid\_nn} model family covers any architecture in which a classical deep network (transformer, recurrent, or convolutional, not limited to LLM-based extractors) performs the main representational work and the quantum circuit is a submodule. A bare variational classifier over precomputed features is \texttt{vqc}, even when those features come from a deep model. Compositional (DisCoCat-style) studies are identified by circuit topology derived from a grammatical parse, regardless of toolkit. These disambiguation rules were fixed during protocol definition and applied uniformly during extraction.

Extraction and classification were performed by one author and reviewed by a second; disagreements were resolved by discussion against the decision rules above. All figures below are reported over the full corpus of 125 studies.

The corpus is small enough that many of the cross-tabulations below have cells in the low single digits, and every count is a complete enumeration of the corpus rather than an estimate from a sample of it. All distributions are therefore reported descriptively, as counts and shares of a known population. Where the article asserts that two dimensions are coupled rather than simply describing one, the claim rests on the size of the imbalance in the cross-tabulation together with a stated mechanism, and both are given. Percentages computed on denominators below ten are reported with their denominators and are not read as evidence of a difference.

\section{Taxonomy of Text Encoding Strategies}
\label{sec:taxonomy}

The encoding problem in QNLP is not a single operation. Depending on the study, text can enter the quantum computation as rotation parameters, as amplitudes of a prepared state, as a grammar-derived circuit, as a structured feature map, or through a more specialized interface such as repeated data injection, a hybrid neural representation, or an observable. These mechanisms are often discussed under the same term, ``encoding,'' although they make different assumptions about where linguistic information is represented and which quantum resource is being used.

This section shows how these mechanisms are organized into the proposed taxonomy as shown in Table~\ref{tab:encoding_taxonomy}, applying the classification rules of Section~\ref{subsec:analysis_framework}: every study discussed below is assigned to exactly one encoding family by the dominant-interface decision rule defined there. The classification is based on the point at which text or text-derived features enter the quantum computation. A study is therefore not classified by the optimizer, measurement rule, downstream classifier, or ansatz alone. These elements are treated as design choices that can be combined with an encoding category.

The taxonomy comprises seven encoding families. Four are major: angle encoding ($n=51$), amplitude encoding ($n=28$), syntactic/DisCoCat encoding ($n=29$), and structured feature-map encoding ($n=9$). Three are emerging families that together account for the remaining eight studies: data re-uploading ($n=4$), hybrid classical--quantum embedding ($n=3$), and Hamiltonian/observable encoding ($n=1$). To use the taxonomy, one may first identify whether the text representation is loaded as a flat vector, imposed through linguistic structure, or used to define a kernel geometry. The next design question is whether the encoding is intended to minimize qubit count, preserve linguistic composition, define a kernel geometry, or improve trainable circuit expressivity. Not all combinations are equally practical on near-term devices. The main tradeoff is between width, depth, state-preparation cost, and the amount of linguistic structure preserved by the quantum circuit \cite{preskill2018quantum,cerezo2021variational,ruskanda_evaluating_2025}.

\begin{table*}[t]
\caption{Taxonomy of text encoding strategies in the final primary corpus: four major families and three emerging families.}
\label{tab:encoding_taxonomy}
\centering
\scriptsize
\begin{tabular}{@{}p{0.18\textwidth} p{0.09\textwidth} p{0.45\textwidth} p{0.20\textwidth}@{}}
\hline
\textbf{Encoding family} & \textbf{Corpus count} & \textbf{Core mechanism} & \textbf{Typical model families} \\
\hline
\multicolumn{4}{@{}l}{\emph{Major families}} \\
Angle encoding & 51 & Text features are mapped to single-qubit rotation parameters, typically through $R_x$, $R_y$, or $R_z$ gates. & VQC, hybrid neural, recurrent neural, quantum transformer \\
Amplitude encoding & 28 & A normalized feature vector is loaded into the amplitudes of an $n$-qubit state in a $2^n$-dimensional Hilbert space. & VQC, quantum kernel, quantum transformer, hybrid neural \\
Syntactic/DisCoCat encoding & 29 & Sentence grammar determines a string diagram or circuit; word meanings are attached to grammatical types and composed by reductions. & Compositional QNLP \\
Structured feature maps & 9 & Classical text features parameterize structured circuits such as IQP, ZZFeatureMap, ZFeatureMap, PauliFeatureMap, or Pauli-Y variants. & Quantum kernel, VQC, hybrid neural, quantum transformer \\
\multicolumn{4}{@{}l}{\emph{Emerging families}} \\
Data re-uploading & 4 & Input features are injected repeatedly across circuit layers, alternating data-dependent rotations with trainable transformations. & VQC, recurrent neural \\
Hybrid embedding & 3 & A classical encoder first produces a compact text representation that is then passed to a quantum layer. & Hybrid neural, VQC \\
Hamiltonian/observable encoding & 1 & Text-derived information is represented through Pauli-string Hamiltonians or observables rather than only through prepared quantum states. & VQC \\
\hline
\end{tabular}
\end{table*}

\subsection{Angle Encoding}
\label{subsec:angle_encoding}

Angle encoding is the dominant category in the corpus. In this scheme, a classical vector $\mathbf{x} = (x_1,\ldots,x_d)$ is mapped to rotation parameters,
\begin{equation}
|\phi(\mathbf{x})\rangle =
\prod_{j=1}^{d'} R_{\alpha_j}(s_j(\mathbf{x})) |0\rangle^{\otimes n},
\end{equation}
where $R_{\alpha_j}$ is a rotation around the $x$, $y$, or $z$ axis and $s_j(\cdot)$ denotes scaling, projection, or feature selection. The input is rarely raw text. It is usually a TF-IDF vector, a word or sentence embedding, a transformer representation, or a compact feature vector produced by a preceding classical model.

The prevalence of angle encoding follows from its implementation simplicity. Rotations fit naturally into parameterized quantum circuits and support the standard gradient-based training workflow of variational algorithms, and they avoid arbitrary-state preparation, which makes them convenient for small simulated or hardware-oriented experiments \cite{mitarai2018quantum,schuld2019evaluating}. What makes the family a single encoding category is this shared rotation mechanism; within it, the corpus studies differ along two encoding-level axes: the rotation construction itself and the classical source that supplies the angles.

The first axis is the rotation construction. The most common pattern assigns one feature to one single-axis rotation, typically an $R_Y$ rotation as in the default angle-embedding routine \cite{ruskanda_evaluating_2025,kesiku_quantum_2025,diaz-ortiz_strongly_2024}, though $R_X$ is also used \cite{day_qnet_2023,zhou_new_2022,tiwari_quantum_2024} and some studies apply combined $R_Y$ and $R_Z$ rotations per qubit \cite{stein_applying_2023}. A denser variant stores more than one feature on each qubit through successive rotations about different axes, trading a shorter register for greater loading depth \cite{kesiku_unlocking_2026}.

The second axis is the classical source of the angles, and it is where the studies vary most, because the rotation parameters are numbers rather than text and so the encoding's content is fixed by the front-end that produces them. Some studies load sparse lexical statistics such as TF-IDF, bag-of-words, or $n$-gram vectors \cite{boopalamani_quantum-enhanced_2025,payares_quantum_2023,kabir_qtopic_2026}; others load static word embeddings \cite{di_sipio_dawn_2022,diaz-ortiz_strongly_2024,ohno_quantum_2025}; and a large, growing group compresses contextual transformer embeddings, for instance reducing a 768-dimensional BERT or DistilBERT vector to a handful of features before rotation \cite{ahmad_hybrid_2023,kumar_hybrid_2026,suleiman_hdqe_2026,yu_application_2024}. A further group derives the angles from deeper classical feature extractors or from handcrafted linguistic features \cite{rezaei_hybrid_2026,mukesh_qvila_2024,mompelat_compact_2026}. In each case a dimensionality-reduction step, whether PCA, a linear projection, or an autoencoder, sits between the text representation and the rotations and is effectively part of the encoding, since it determines what survives into the quantum circuit.

The limitation is inherent in this design. Encoding cost grows with the number of retained features, so a model must spend more rotations, more qubits, or more layers as it keeps more of the input. Angle encoding therefore pushes the hard part of the problem onto classical dimensionality reduction, which makes it practical for NISQ-scale prototypes but also means the quantum circuit often operates on a heavily compressed view of the linguistic input.


\subsection{Amplitude Encoding}
\label{subsec:amplitude_encoding}

Amplitude encoding maps a normalized vector $\mathbf{x} \in \mathbb{R}^{2^n}$ or $\mathbb{C}^{2^n}$ to a quantum state,
\begin{equation}
|\phi(\mathbf{x})\rangle = \sum_{i=0}^{2^n-1} x_i |i\rangle,
\quad \sum_i |x_i|^2 = 1.
\end{equation}
For NLP, the appeal is immediate. A high-dimensional sentence embedding can, in principle, be represented with logarithmically many qubits: a 768-dimensional transformer vector, for example, can be padded and normalized into a ten-qubit state. This density is why amplitude encoding is most often paired with dense contextual representations, above all BERT, ELECTRA, and other transformer embeddings, and to a lesser extent with static word embeddings or TF-IDF and bag-of-words vectors \cite{buonaiuto_quantum_2024,guarasci_applying_2025,cavar_old_2026,zeguendry_quantum_2024}.

Because the register holds exactly $2^n$ amplitudes, the encoding is defined as much by how the source vector is fitted to that dimension as by the loading step itself. When the embedding is already close to a power of two, it is simply normalized and zero-padded, as when a 768-dimensional vector is loaded into ten qubits \cite{guarasci_applying_2025,buonaiuto_multilingual_2025}. When it is not, or when the qubit budget is smaller, a classical reduction such as PCA, t-SNE, or UMAP is applied first, so that the reduction and the amplitude map together fix the geometry the circuit actually sees \cite{korrapati_quantum_2025,zhang_entangled_2024,tufchi_quantum-enhanced_2026}.

The reason to accept this encoding is representational. Storing a large vector in a small register makes overlap- and inner-product-based operations natural, because the normalized state \emph{is} the embedding. SWAP-test similarity and quantum-kernel pipelines exploit this directly, letting the geometry of the encoded state stand in for the similarity structure of the embedding space \cite{cavar_word_2025,zeguendry_quantum_2024,schuld2019quantum}. A distinct thread goes further and encodes complex-valued rather than real-valued amplitudes, representing tokens as complex state vectors and building complex-valued semantic operations on top of them, which uses a property of the amplitude space that has no direct classical-vector analogue \cite{routh_quantum_2026,cavar_old_2026}.

The signature difficulty is state preparation, and it is where the studies differ most in how honestly the cost is treated. Most rely on exact library routines that prepare an arbitrary amplitude vector, with the preparation depth absorbed into a simulator call \cite{buonaiuto_quantum_2024,guarasci_applying_2025}. A smaller set confronts the cost directly through divide-and-conquer amplitude encoding, which trades circuit depth against ancilla width and is validated on \texttt{ibm\_kyoto} for a word-sense disambiguation subset alongside the noiseless-simulation classification results \cite{zhang_entangled_2024}. Loading an arbitrary vector generally requires a circuit whose depth can offset the qubit savings, so experiments that hide this behind high-level routines remain useful for comparing representation strategies under a fixed qubit budget but should not be read as hardware-ready evidence unless the preparation cost is included.

\subsection{Syntactic and DisCoCat Encoding}
\label{subsec:syntactic_encoding}

Syntactic encoding is the characteristic interface of compositional QNLP. Instead of first reducing a sentence to a flat feature vector, the model begins with a grammatical derivation. Words are assigned grammatical types, the sentence is reduced according to a pregroup or related grammar, and the resulting string diagram is translated into a tensor network or quantum circuit \cite{coecke2010mathematical}. What makes the family a single encoding category is a shared two-stage construction: a diagram that fixes the circuit \emph{topology} from grammar, followed by a rule that turns that diagram into a parameterized qubit circuit. The 29 studies in the corpus differ mainly along these two stages, and it is these differences, rather than the downstream task, that define the encoding-level variants below.

The first stage is how the diagram is obtained. The canonical route is a pregroup derivation, in which each word receives a pregroup type and grammatical reductions become the wires that connect word boxes; pregroup parsing is efficient and underlies most DisCoCat implementations \cite{meichanetzidis_quantum_2021,meichanetzidis_grammar-aware_2023,lorenz_qnlp_2023}. In practice this parse is often produced by a broader-coverage CCG parser and then rewritten into DisCoCat form, as in studies that derive their diagrams with the Bobcat or DepCCG parsers before compilation \cite{abiwardani_generalized_2025,ruskanda_quantum_2022,r_quantum_2023}. A second variant replaces the parse with a fixed \emph{reader} that imposes a uniform composition rule instead of syntax: a bag-of-words reader composes words through a single commutative spider, and a word-sequence reader composes them left to right into a tensor train, both of which appear as structure-free counterparts to the grammar-based diagrams \cite{lorenz_qnlp_2023,liu_research_2024,peral-garcia_comparing_2024}. These readers trade syntactic fidelity for shallower, more uniform circuits and are frequently used as structure-free reference points within DisCoCat studies. A third variant changes the scope of the diagram itself, lifting the sentence-level construction to the discourse level so that pronouns and cross-sentence reference are represented, as in DisCoCirc-style coreference models \cite{wazni_large_2024}.

The second stage is how the diagram becomes a qubit circuit. Each atomic type is assigned a fixed number of qubits, so the noun and sentence types set the circuit width, grammatical reductions compile to Bell-style effects between wires, and function words such as relative or possessive pronouns are realized as Kronecker or GHZ spiders \cite{meichanetzidis_grammar-aware_2023}. Word boxes are then filled by a parameterized ansatz, and the dominant choices are the IQP ansatz used by default in \textsc{lambeq} and, where circuit expressivity is compared, the Sim-14, Sim-15, and strongly entangling ans\"atze \cite{lorenz_qnlp_2023,castillo_comparative_2025}. Some studies design task-specific ans\"atze on top of the same grammatical skeleton rather than proposing a new diagram \cite{ruskanda_simple_2023,abiwardani_generalized_2025}. A final variant intervenes between the two stages by rewriting the diagram before compilation: ZX-calculus is used to simplify the grammar-derived circuit and reduce its resource cost \cite{sreedhar_towards_2025}.

This category differs from vector-loading schemes. The circuit is not only a host for numerical features; its wiring reflects syntactic composition. This gives syntactic encoding a clear interpretability advantage. Grammatical roles, reductions, and word interactions remain visible in the model structure, which explains why nearly all compositional papers in the corpus use this category.

The tradeoff is flexibility. Syntactic encoding depends on parsers, grammar formalisms, and sentence structures that can be translated into manageable circuits. Many experiments therefore use short, controlled, or grammar-friendly datasets. As sentences become longer, noisier, multilingual, or document-level, the parser and circuit-growth requirements become more restrictive. Thus, syntactic encoding is the most linguistically structured category, but it is also the category that most depends on the availability of a reliable grammar-to-circuit pipeline.

\subsection{Structured Feature-Map Encoding}
\label{subsec:iqp_encoding}

Structured feature maps encode data through circuits with phase rotations, Pauli operators, and entangling gates whose parameters depend on products or interactions of input features, so that the encoded state defines an implicit feature space in which examples are compared by state similarity \cite{schuld2019quantum}. Within the corpus the map is drawn from a small set of families. The most common is the ZZ feature map, which introduces pairwise feature interactions through $ZZ$ entangling rotations \cite{widdows_near-term_2024,baik_neural_2025,sunkavalli_experimentation_2024}. IQP embeddings, which apply diagonal phase gates in the Hadamard basis, are used similarly \cite{djemmal_arabert-qc_2025,ruskanda_quantum-enhanced_2023}. A few studies construct higher-order maps from selected Pauli terms, for example combined $X$, $Y$, and $ZZ$ gates or a hierarchy of Pauli-$Y$ and second-order $YY$ interactions, in order to encode richer correlations \cite{katyayan_supervised_2023,ruskanda_quantum-enhanced_2023}. One study instead reads the feature map out through local observables, using a projected quantum kernel to obtain a lower-variance similarity estimate \cite{keith-norambuena_quantum_2026}.

Because both the qubit count and the number of encoded interactions grow with the input dimension, a structured feature map is almost always preceded by an aggressive classical reduction. In practice this is usually PCA to on the order of ten to a few dozen components, and occasionally a wavelet or related projection, applied to an already compressed text representation such as a TF-IDF vector or a reduced transformer embedding \cite{mohammad_masum_hybrid_2023,djemmal_arabert-qc_2025,jagatheesan_quantum_2026}. The reduction is thus part of the encoding, since it fixes both the qubit budget and which feature interactions the map can express.

The benefit is that feature interactions are built into the map. Rather than using the circuit only to host features, the feature map defines the geometry in which examples are compared, which makes the approach attractive for testing whether a quantum kernel yields a useful similarity space for text.

The limitation is that the feature map becomes a strong modeling assumption. Performance depends on whether the chosen interactions are relevant to the linguistic task and on the preceding compression, so a structured feature map offers a principled kernel-oriented interface whose usefulness is set jointly by map design and the reduction that feeds it.

\subsection{Emerging Encoding Families}
\label{subsec:other_encoding_interfaces}

The three remaining families are individually too small to rank among the major families; the taxonomy classifies them as emerging, each defined by a distinct mechanism. Data re-uploading addresses the small-register setting by injecting the input features multiple times across the depth of a circuit: rather than loading the data once and then applying only trainable gates, the model alternates data-dependent rotations with trainable transformations, so expressivity grows with depth rather than with qubit count. The four corpus studies differ mainly in what is re-uploaded and where. One implements a bricklayer of dense rotation gates in which repeated feature injection, rather than one-shot angle loading, is the defining encoding choice \cite{hamilton_characterizing_2023}. A recurrent variant re-uploads code-token embeddings across successive quantum embedding layers, so the same feature stream re-enters the circuit at every recurrent step \cite{song_recurrent_2024}. A third injects lexical features repeatedly into a seven-qubit NISQ classifier, explicitly trading qubit width for circuit depth \cite{silver_lexiql_2024}, and a fourth re-uploads sentence-transformer embeddings across the layers of an eighteen-qubit circuit, using qubit count and depth jointly to control capacity in a quantum LLM fine-tuning setting \cite{kim_quantum_2025}. The shared advantage is that expressivity can be raised without more qubits; the shared drawback is that the encoding burden moves from width to depth.

Hybrid embedding is a system-level interface rather than a single gate-level rule. A classical model first produces a compact representation, and a quantum layer then consumes it through an angle, amplitude, or less explicitly specified feature-injection mechanism. In the corpus the three studies are distinguished mainly by how aggressively the classical stage compresses and how clearly the quantum loading is described. One compresses fused text and image embeddings into an eight-dimensional latent vector and matches it to an eight-qubit VQC for emotion recognition in dialogue \cite{castrignano_toward_2025}. Another compresses GloVe and fastText code-token embeddings into a minimal intermediate representation before a four-qubit vulnerability classifier, without specifying the gate-level loading rule \cite{akter_automated_2023}, and a third fuses GloVe text features with VGG16 visual features through a quantum layer for misinformation detection, again with no stated state-preparation method \cite{bikku_quantum-enhanced_2025}. This underspecification is characteristic of the family and feeds its main difficulty, attribution: when most semantic processing happens in the classical stage, the independent contribution of the quantum layer is hard to isolate.

Hamiltonian and observable encoding is even more specialized. The single corpus study maps text-derived vectors onto Pauli-string Hamiltonians and reads them out as quantum observables rather than preparing them as amplitude or rotation states, which avoids conventional state encoding entirely and is reported to keep qubit and gate complexity logarithmic \cite{tiblias_efficient_2026}. This relocates the loading problem from state preparation to observable design and may sidestep some state-preparation costs, but with only one study in the corpus the evidence base is too thin to characterize the family beyond this single design point.

\subsection{Taxonomic Implications}

The taxonomy makes clear that encoding is not a neutral preprocessing choice in QNLP. It determines which linguistic information reaches the quantum circuit, how much classical compression is required, and where the main resource bottleneck appears. Angle and amplitude encodings start from numerical text representations and differ mainly in how they trade circuit width, depth, and state-preparation cost. Syntactic encoding instead gives the grammatical analysis a structural role in the circuit. Structured feature maps define the comparison space used by kernel methods, while the emerging families modify these patterns through repeated data injection, classical--quantum embedding layers, or observable-based representations.

This distinction matters because many cross-paper comparisons implicitly compare different views of the input rather than only different learning architectures. A VQC using compressed word or sentence vectors, a DisCoCat circuit generated from a parser, and a kernel model built around a designed feature map make different commitments about what should be preserved from the original text. Their results are therefore shaped jointly by the classical representation, the compression method, the encoding circuit, the ansatz, the measurement rule, and the execution backend.

The taxonomy is therefore used in the remainder of the survey to organize the evidence, not to rank the encodings. The available studies do not yet show that any single family is generally superior across QNLP tasks. Instead, they show that the encoding stage frames the experiment itself: it defines the form of the linguistic input, constrains the feasible circuit resources, and conditions how performance and scalability results should be read.

\Figure[t!](topskip=0pt, botskip=0pt, midskip=0pt)[width=\textwidth]{figures/encoding_by_model_family.pdf}{Encoding family by model family in the final primary corpus. Cell values report the number of studies assigned to each encoding--model-family pair.\label{fig:encoding_model_family}}

\section{Comparative Analysis of Encoding Families}
\label{sec:comparative_analysis}

This section uses the taxonomy to assess how far the main encoding families have developed in practice. Rather than revisiting the definition of each family, the focus is on what the corpus reveals about their use: which encodings have become standard design choices, which remain exploratory, and how strongly the reported results are supported by datasets, baselines, hardware assumptions, and scalability analysis.

The analysis applies the classification scheme of Section~\ref{subsec:analysis_framework}. Each study is assigned to its dominant encoding interface, although most implementations combine several steps before measurement: text representation, feature reduction, circuit encoding, a trainable or fixed quantum model, and a readout rule. This matters because an observed result is rarely attributable to the encoding alone. It often reflects the full pipeline in which the encoding is embedded.

Before examining the families individually, this section establishes that the choice of encoding is not independent of the model family in which it appears. As shown in Figure~\ref{fig:encoding_model_family}, which crosses the two dimensions, the resulting distribution is far from uniform: most encodings concentrate in one or two model families rather than spreading evenly across them. The three dominant encodings account for 108 of the 125 studies, and the same coupling holds among them, though only one of the three, syntactic/DisCoCat encoding, is fixed to its architecture as an artifact of the classification rather than as a finding about the field: it and the compositional model family coincide exactly, all 29 studies of the one being the same 29 studies as the other, because Section~\ref{subsec:analysis_framework} identifies both categories by the same property, a circuit topology derived from a grammatical parse, so the two labels cannot come apart by construction. The coupling claim for the remaining two dominant encodings therefore has to rest on the rest of the corpus, and it does, though less absolutely. Excluding the compositional studies, the coupling persists across the remaining 96: every encoding family except angle places a majority of its studies in a single model family, on a subset where VQCs alone absorb 42 of the 96. Structured feature maps tie to quantum kernels, where six of the nine structured-map studies live. Angle encoding, by contrast, is the general-purpose default shared across the remaining gate-model architectures, and amplitude encoding co-resides with it inside VQCs and kernel methods. Comparing ``encoding families'' across the corpus is therefore partly comparing the model families they are bound to. This motivates the approach of the four subsections that follow: each holds one established family fixed and asks how its use changes across the contexts in which it appears. The emerging families are treated last.

\subsection{Angle Encoding Across Model Families}
\label{subsec:angle_across_families}

Angle encoding is the one interface that is not bound to a single architecture, so it is the natural place to begin: rather than asking which encoding each model family prefers, the analysis holds the encoding fixed and asks how the same data-dependent rotation is repurposed as the surrounding model changes. The 51 angle-encoding studies distribute across five model families, VQCs (20), hybrid neural models (10), recurrent models (10), quantum transformers (9), and quantum kernels (2), and their treatment of the rotation differs enough that ``angle encoding'' names a role rather than a fixed operation. As shown in Table~\ref{tab:angle_by_family}, which summarizes the split, the remainder of this subsection reads it family by family. Qubit budgets are reported by 34 of the 51 studies and, with one structural exception, stay within a narrow $2$--$8$ band, so the differences below are differences in \emph{where} the rotation sits, not in how wide a register it drives.

\begin{table*}[t]
\caption{Angle encoding split by model family in the final primary corpus.}
\label{tab:angle_by_family}
\centering
\scriptsize
\begin{tabular}{@{}p{0.15\textwidth} p{0.06\textwidth} p{0.13\textwidth} p{0.26\textwidth} p{0.27\textwidth}@{}}
\hline
\textbf{Model family} & \textbf{Studies} & \textbf{Median qubits (range)} & \textbf{Typical classical front-end} & \textbf{Role of the angle rotation} \\
\hline
VQC & 20 & 4.5 (2--16) & TF-IDF, BERT/SBERT/DistilBERT, handcrafted features & Sole quantum front door before a trainable ansatz \\
Hybrid neural & 10 & 4.5 (2--8) & Output of a classical transformer/RNN encoder & Submodule consuming an already-compact representation \\
Recurrent neural & 10 & 4 (4--8) & Word or BERT embeddings, per token & Data input inside the recurrent cell, re-entered each step \\
Quantum transformer & 9 & 4 (4--28) & Word embeddings, multimodal features & State preparation for quantum attention \\
Quantum kernel & 2 & 4 (4) & TF-IDF or compressed feature vectors & Definition of a similarity-circuit geometry \\
\hline
\end{tabular}
\end{table*}

In VQCs, angle encoding appears in its canonical form: the rotation layer \emph{is} the quantum front door, followed only by a trainable ansatz and a measurement, so the encoding carries the entire burden of moving text into the circuit. This is also the family with the widest spread of classical front-ends, which is why the compression the encoding demands is most visible here. Studies load TF-IDF or bag-of-words features \cite{thota_comparison_2025}, transformer embeddings reduced by one or two orders of magnitude, from BERT's 768 dimensions to ten \cite{ahmad_hybrid_2023}, from DistilBERT to four \cite{kumar_hybrid_2026}, or as far down as two values before an $R_x$ layer \cite{altintas_beyond_2025}, and in a few cases handcrafted linguistic descriptors such as character entropy and diacritic counts \cite{mompelat_compact_2026}. Because the circuit is the model rather than a component of it, the VQC setting also hosts almost all of the corpus's controlled encoding comparisons: a QCNN study finds angle encoding the best of angle, amplitude, and IQP maps \cite{ouamane_sentiment_2025}, and a sentiment study finds an $R_y$ angle map beating a Pauli feature map under matched circuits \cite{ruskanda_evaluating_2025}. The same canonical pattern recurs across further tasks and domains in the corpus. Examples include a dressed-quantum-circuit transfer-learning classifier built on frozen BERT embeddings \cite{ardeshir-larijani_hybrid_2024}; SBERT embeddings fused with Pauli expectation values for scholarly AI-text detection \cite{hazim_hqml-nlp_2026}; angle-encoded word-sense disambiguation on WordNet sentence-sense pairs \cite{karanath_quantum-enhanced_2025}; and a privacy-oriented quantum language model trained on NISQ hardware and transferred to classical downstream tasks \cite{li_pqlm_2023}. Further instances are post-hoc calibration of language-model confidence via a variational circuit \cite{neto_quantum_2024}; a quantum depthwise-convolutional classifier that pairs amplitude-encoded word embeddings with an angle-encoded sentence representation \cite{chen_multi-scale_2025}; noise-robust multimodal sentiment analysis trained with an SPSA optimizer \cite{hussain_quantum_2025-1}; and an all-quantum LSTM with a self-attention block whose noise robustness is tested explicitly \cite{yan_exploring_2026}. The attribution of any result is comparatively clean here, but at the price of a circuit that only ever sees a heavily compressed view of the input.

Hybrid neural models move the rotation downstream. The angle layer is no longer the front door but a submodule embedded in a larger classical network, so a transformer or recurrent encoder performs most of the representational work and the quantum circuit consumes a representation that has already been compressed classically. This is visible both in the front-ends, MARBERT, AraBERT, or XLM-R embeddings compressed into shallow circuits whose Pauli readouts are fused back into the classical model \cite{suleiman_hdqe_2026}, a BERT-and-GRU multimodal pipeline projected into a two-qubit circuit, a transformer generator whose tokens are injected as rotations \cite{routh_hybrid_2025}, and in the registers, which are the smallest in the angle subcorpus. The same downstream-submodule role recurs in further hybrid designs. These include a quantum fuzzy neural network layered on top of angle-encoded features for noise-robust sentiment analysis \cite{dave_sentiqnf_2026}; quantum tensor-network fusion of text, audio, and visual streams for multimodal sentiment \cite{li_qd-msa_2025,li_qmlsc_2025}; and a compact variational embedding layer used as a drop-in replacement for a classical word-embedding layer \cite{roshan_ahmed_deepket-_2024}. Two further designs are a quantum adaptation of the Mamba selective state-space architecture for language modeling \cite{stenzel_qmamba_2025} and a self-supervised, GLUE-benchmarked quantum-classical BERT variant pre-trained under an explicit noise model \cite{yao_self-supervised_2025}. The cost is the mirror image of the VQC case: because the classical encoder dominates, the independent contribution of the quantum rotation is the hardest to isolate of any family.

Recurrent models repurpose the rotation temporally rather than spatially. Here angle encoding is applied \emph{inside} a quantum recurrent cell, a QLSTM or QRNN block, so the same rotation interface is re-entered at every step of the sequence rather than once at the input \cite{stein_applying_2023,kesiku_quantum_2025,yu_application_2024,diaz-ortiz_strongly_2024}. This keeps the register deliberately narrow, a median of four qubits and never more than eight, because width is fixed by the cell while information enters through repeated application along the sequence, which brings this usage closer in spirit to data re-uploading than to a single loading step. The $R_y$-based \textsc{AngleEmbedding} is the recurring choice \cite{kesiku_quantum_2025}, and the family is the one most oriented toward token-level tasks such as part-of-speech tagging \cite{pandey_quantum_2026} alongside sequential sentiment classification, code-mixed part-of-speech tagging across three Indian language pairs \cite{pandey_quantum_2023}, recency-weighted mental-health-aware sentiment monitoring \cite{rawat_mental_2026}, frozen-BERT emotion recognition \cite{vm_hybrid_2026}, and movie-review sentiment classification competitive with classical RNN, GRU, and LSTM baselines \cite{xu_quantum_2024}.

Quantum transformers use the rotation to prepare the states that attention then compares, and this is the one place where the usual qubit discipline breaks. Features are mapped into per-token quantum states whose overlaps play the role of attention scores, so register width scales with sequence and embedding size rather than with a compressed feature count. QNet assigns one qubit per token--embedding-dimension pair and splits axes by role, with $R_x$ carrying the embedding value and $R_z$ a positional phase \cite{day_qnet_2023}. A near-term self-attention study reaches 28 qubits, the widest register in the entire angle subcorpus \cite{he_training_2024}. Other studies keep the per-token construction but vary what the attention operates over: a quantum self-attention network built on the QSANN architecture \cite{li_quantum_2024}, a quantum graph transformer that folds self-attention into a graph message-passing framework \cite{aktar_quantum_2025}, a mixed-density-matrix formulation of quantum self-attention \cite{chen_quantum_2025}, and a noisy-simulation image-text quantum transformer that fuses a text PQC with a vision-transformer PQC \cite{wu_multimodal_2025}. Most implementations avoid this blow-up by keeping the quantum attention block small \cite{di_sipio_dawn_2022,pham_fitqt_2026}, but the family is nonetheless the only one whose structural logic pushes angle encoding beyond the near-universal $4$--$8$ qubit band.

\begin{table*}[t!]
\caption{Amplitude encoding split by model family in the final primary corpus.}
\label{tab:amplitude_by_family}
\centering
\scriptsize
\begin{tabular}{@{}p{0.15\textwidth} p{0.06\textwidth} p{0.13\textwidth} p{0.26\textwidth} p{0.27\textwidth}@{}}
\hline
\textbf{Model family} & \textbf{Studies} & \textbf{Median qubits (range)} & \textbf{Typical classical front-end} & \textbf{Role of the amplitude state} \\
\hline
VQC & 15 & 10 (9--16) & Normalized BERT/ELECTRA/DistilBERT/SBERT embeddings & Input state prepared for a trainable ansatz \\
Quantum kernel & 5 & 10.5 (7--12) & Word/sentence embeddings, TF-IDF vectors & Overlap geometry read out by a SWAP test \\
Quantum transformer & 4 & 7.5 (4--12) & Word or contextual token embeddings & State preparation for quantum attention \\
Hybrid neural & 3 & 9 (6--12) & Multimodal or token embeddings & High-fidelity submodule after a classical encoder \\
Recurrent neural & 1 & --- & Chinese word embeddings & Per-gate state inside a decomposed QLSTM \\
\hline
\end{tabular}
\end{table*}

Quantum kernels are represented by only two angle studies, and the contrast with the previous families is instructive. When angle encoding does appear in a kernel method, the rotation defines a similarity-circuit geometry rather than feeding a trainable ansatz \cite{hsu_quantum_2025,boopalamani_quantum-enhanced_2025}; but kernel methods overwhelmingly prefer purpose-built structured feature maps (Section~\ref{subsec:structured_across_constructions}), so angle encoding is a marginal, borrowed choice here rather than a native one. This is the exception that confirms the pattern: angle encoding travels everywhere, but in the one family whose encoding is itself the model, a different interface wins.

Taken together, the split shows that holding the encoding fixed does not hold the pipeline fixed. The same rotation is a complete front door in VQCs, a downstream submodule in hybrid models, a per-step cell input in recurrent models, an attention-state preparation in quantum transformers, and a borrowed similarity map in kernels. These roles differ in exactly the dimensions that matter for evaluation: how much of the model the encoding accounts for, whether register width is set by compression or by architecture, and how cleanly the quantum contribution can be attributed. The one invariant across all five families is the heavy classical compression that precedes the circuit. This confirms the encoding--model-family correlation of Fig.~\ref{fig:encoding_model_family} from the opposite direction: even when the encoding is the same, it is the model family that determines what the encoding does, which is why angle-encoded results are not directly comparable across architectures despite sharing a label.

\subsection{Amplitude Encoding Across Model Families}
\label{subsec:amplitude_across_families}

Amplitude encoding is the natural counterpoint to angle encoding, and reading its 28 studies through the same per-family lens makes the contrast concrete. It appears in the same host families, VQCs (15), quantum kernels (5), quantum transformers (4), hybrid neural models (3), and a single recurrent model, but with an almost inverted profile, both in which families it favors and in the resources it consumes. Where angle encoding compresses text to a handful of rotation parameters over a $2$--$8$ qubit register, amplitude encoding loads a full normalized vector into the $2^{n}$ amplitudes of an $n$-qubit state, so its registers cluster near $\log_2$ of the embedding dimension rather than at a fixed small width. As shown in Table~\ref{tab:amplitude_by_family}, which summarizes the split, the immediate signal is the qubit column: every family sits at a median near ten qubits, against the median of four seen for angle encoding, because the encoding is chosen precisely to avoid the aggressive projection that angle encoding requires.

VQCs are again the dominant host, but the role differs from the angle case. Amplitude encoding here prepares the input state that a trainable ansatz, typically \textsc{EfficientSU2} or \textsc{RealAmplitudes}, then classifies \cite{korrapati_quantum_2025}. The front-end is almost exclusively a dense contextual embedding normalized and padded to the nearest power of two: a 768-dimensional SBERT vector into ten qubits \cite{guarasci_applying_2025}, BERT or ELECTRA embeddings padded from 768 to 1024 amplitudes over ten qubits \cite{buonaiuto_quantum_2024}, and mBERT embeddings reused across multilingual transfer \cite{buonaiuto_multilingual_2025}. The $\log_2$ packing is stated as the point of the method rather than an implementation detail \cite{liu_hcqnsc_2026}. Where a study compares the two interfaces directly, it frames amplitude encoding as the high-fidelity option and angle encoding as the cheaper fallback: one sentence-classification model reports its amplitude variant at $0.81$ accuracy against $0.71$ for an MLP-reduced angle variant of the same system \cite{guarasci_applying_2025}. The same host role recurs across further tasks. Examples include an attention-augmented amplitude classifier evaluated on a non-standard Kaggle text corpus \cite{amsaveni_qesan_2025}; an Electra-based amplitude transfer-learning system for Italian-language sentiment analysis \cite{buonaiuto_quantum_2024-1}; and an amplitude-encoded QRNN with ancilla-mediated state transfer applied to code-mixed social-media text \cite{pandey_quantum_2024}. Further instances are a complex-amplitude variational topic model that reduces parameter count relative to a classical neural topic model while improving topic coherence \cite{qiao_quantum_2024}; a BERT-to-nine-qubit amplitude pipeline for smishing and spam-text detection \cite{shahriyar_toward_2025}; and a ten-qubit amplitude classifier applied to word-sense disambiguation in a low-resource ancient language \cite{isea_using_2026}. This family is also markedly more recent than its angle counterpart, with every VQC amplitude study appearing from 2024 onward, which tracks the point at which transformer embeddings became the default text representation and made a dense-vector encoding worth its qubit cost.

Quantum kernels are where amplitude encoding is most at home, in contrast to the marginal, borrowed status angle encoding held there. When a normalized embedding is amplitude-encoded, the inner product of two states equals the inner product of the underlying vectors, so a SWAP test or state-overlap estimate reads out cosine-like semantic similarity directly, with no training involved \cite{cavar_old_2026,cavar_word_2025,cavar_semantic_2024}. This is why the semantic-similarity task cluster of the corpus lives almost entirely in this cell, and why kernel studies apply the same construction to fastText, GloVe, or BERT embeddings and to TF-IDF vectors feeding a quantum nearest-neighbor classifier \cite{alexander_quantum_2022,zeguendry_quantum_2024}. The fit is principled rather than incidental: the encoded state \emph{is} the geometry the kernel measures, and some studies push the density further with complex-amplitude variants that map pairs of dimensions to the real and imaginary parts of each amplitude, halving the register again \cite{cavar_old_2026}.

\begin{table*}[t!]
\caption{Syntactic/DisCoCat encoding split by circuit ansatz in the final primary corpus.}
\label{tab:syntactic_by_ansatz}
\centering
\scriptsize
\begin{tabular}{@{}p{0.20\textwidth} p{0.06\textwidth} p{0.13\textwidth} p{0.30\textwidth} p{0.21\textwidth}@{}}
\hline
\textbf{Ansatz family} & \textbf{Studies} & \textbf{Median qubits (range)} & \textbf{Gate-level construction} & \textbf{Noun parameterization} \\
\hline
IQP (\textsc{lambeq} default) & 19 & 7 (2--25) & Hadamard layers alternated with layers of parameterized controlled-$R_z$/CZ rotations, one block per grammatical reduction & Nouns are trainable rotations, same as other word types \\
Named alternative (SimpleSA, GeneralQC, DCA) & 3 & 5 (4--10) & $H$--CNOT--$R_z$--$H$ or $H$--$CR_x$/$R_x$--CNOT/$R_x$--$CR_z$ blocks & Nouns fixed as identity gates or frozen embedding-derived states \\
Multi-ansatz benchmark (Sim14/15, StronglyEntangling, MPS/Spider/Tensor vs. IQP) & 5 & 13 (7--20) & Same diagram compiled through several interchangeable circuit or tensor-network templates, compared head-to-head & Nouns trainable in every compared variant \\
ZX-calculus/underspecified & 2 & 10 (single report) & Diagram or grammar tree compiled to a circuit via a generic, unlabeled routine; gate template not stated & Not reported \\
\hline
\end{tabular}
\end{table*}

Quantum transformers use amplitude encoding, as they use angle encoding, to prepare the states that attention compares, but over noticeably smaller and more variable registers than in the VQC case \cite{zheng_design_2023,sharma_elevating_2025,aktar_hybrid_2026,karanath_q-optbert_2026,karanath_q-entangleformer_2026}. Amplitude encoding is rarer here than angle encoding (four studies against nine) because preparing an arbitrary state for every token is more expensive than injecting a rotation. The studies that do adopt it tend to encode a single pooled or sentence-level vector rather than a per-token sequence. Hybrid neural models place amplitude encoding in the same downstream-submodule role that angle encoding occupies there, but selected explicitly for fidelity. One knowledge-graph study contrasts an amplitude model chosen for high-dimensional fidelity against an angle model chosen for gate efficiency, the trade-off between the two families stated within a single system \cite{hendl_incorporating_2025}. A multimodal model amplitude-encodes fused visual and textual features \cite{zheng_quantum_2025}, and a space-efficient hybrid architecture amplitude-encodes representations for natural-language-to-Prolog query translation, discussing the state-preparation cost explicitly as a limitation \cite{ahmed_quantum_2024}.

The recurrent family is the sharpest point of the contrast. Amplitude encoding appears in exactly one recurrent study, a decomposed quantum LSTM over Chinese word embeddings \cite{hou_realization_2022}, against the ten recurrent studies that use angle encoding, and the asymmetry is not accidental: re-preparing an arbitrary amplitude state at every step of a sequence is prohibitively expensive, whereas re-applying a rotation inside a cell is cheap. The encoding that is native to recurrent processing is therefore the opposite of the one that is native to overlap-based comparison. Read against the angle-encoding split, amplitude encoding partitions the same families by an opposite logic: angle encoding minimizes width and pays for it in classical projection, while amplitude encoding minimizes projection and pays for it in register width and, more importantly, in state-preparation depth. That last cost is the one most often abstracted away by simulator calls, so amplitude-encoded results are the ones whose reported qubit efficiency is most sensitive to whether the price of preparing the state has actually been counted.

\subsection{Syntactic/DisCoCat Encoding Across Ansätze}
\label{subsec:syntactic_across_ansatze}

Syntactic/DisCoCat encoding cannot be split by model family the way the previous two families were, because all 29 of its studies belong to a single model family: compositional QNLP. This follows from the definition given earlier, since the grammatical reduction that wires the circuit is the model, so there is no separate architecture to vary. What does vary, and what plays the structural role model family played above, is the ansatz used to turn each grammatical box into gates, because that choice is the one genuine degree of freedom left once the parser has fixed the diagram. Two properties are shared across the whole subcorpus regardless of ansatz. First, the toolchain is a near-monoculture: 22 of the 29 studies build their circuits with \textsc{lambeq}, and most of the rest use the related \textsc{DisCoPy}/\textsc{pytket} stack, in contrast to the diverse PennyLane, Qiskit, and TorchQuantum backends seen for the vector-loading families. Second, qubit count is not a function of the ansatz: it is fixed by the diagram, growing with sentence length and the number of grammatical wire types regardless of which gate template is dropped onto each box. As shown in Table~\ref{tab:syntactic_by_ansatz}, which groups the studies by ansatz, the qubit column shows exactly this: the reported range is nearly identical in the IQP bucket and in the small named-alternative bucket, two to twenty-five against four to ten, because register width tracks the sentence, not the gate design.

The IQP ansatz is the family's workhorse and accounts for two-thirds of the subcorpus. It originates in the two ``generic families'' of ansätze proposed alongside the first NISQ demonstrations: a CNOT-and-$U(3)$ construction and an IQP construction of alternating Hadamard and controlled-phase layers \cite{meichanetzidis_quantum_2021}. The IQP variant was then formalized as a depth-$d$ word-circuit template (arity-one words as a single Euler rotation, higher-arity words as $d$ layers of Hadamard-plus-neighboring-controlled-$R_z$ gates), motivated explicitly by IQP's conjectured classical hardness \cite{meichanetzidis_grammar-aware_2023}. It was tuned further with a cheaper single-rotation option for low-arity words \cite{lorenz_qnlp_2023}. This is the construction that \textsc{lambeq} packages as its default \texttt{IQPAnsatz} class \cite{kartsaklis_lambeq_2021}. The great majority of application papers in the corpus adopt it unmodified, changing only what is fed into the diagram or read out of the circuit. Examples include video-caption classification \cite{r_quantum_2023}; multimodal captioning with a separate image encoder \cite{hawashin_compositional_2025}; knowledge-graph relation modeling \cite{bharti_quantum_2025}; sentiment classification that keeps the same \texttt{IQPAnsatz} for its quantum runs while using a Spider tensor-network ansatz only as a classical simulation baseline \cite{ganguly_quantum_2022}; Hindi sentiment classification \cite{naik_extending_2026}; and natural-language inference with tied parameters \cite{sun_quantum_2026}. Further instances are optimizer benchmarking at fixed ansatz depth \cite{rajashekharaiah_sentence_2023,ruskanda_enhancing_2024}; an ensemble over five diagram-construction readers that all compile to the same IQP circuit \cite{bouakba_ensemble_2023}; and a hand-built IQP-equivalent circuit for pronoun resolution built from Fock-space projections \cite{wazni_quantum_2022,wazni_large_2024}. The construction also travels across languages and modalities: a direct translation of the same block structure to Chinese grammar \cite{balthasar_mansky_adapting_2023}, an English--Persian translation pipeline that Euler-decomposes noun rotations \cite{abbaszadeh_hybrid_2026}, and a musical-phrase model that applies the same Euler decomposition to snippet ``words'' \cite{miranda_quantum_2022}. The one recurring modification within this bucket is additive rather than structural, such as a dedicated negation gadget layered on top of an otherwise standard IQP circuit \cite{ruskanda_quantum_2022}, which confirms that once fixed by the toolkit, the ansatz itself is treated as infrastructure rather than as a variable most studies think to manipulate.

A second, much smaller cluster treats the ansatz as the thing to redesign, and does so along one consistent axis: removing the noun from the pool of trainable gates. SimpleSA replaces the IQP block with an $H$--CNOT--$R_z$--$H$ construction in which noun boxes are fixed identity gates rather than trainable unitaries, and reports higher accuracy than IQP with fewer parameters \cite{ruskanda_simple_2023}. GeneralQC generalizes the idea to three concrete gate layouts: $H$-$CR_x$, $R_x$-CNOT, and $R_x$-$CR_z$. The two best-performing variants, $R_x$-CNOT and $R_x$-$CR_z$, are again the ones that leave nouns non-parameterized, reaching 96.91\% accuracy against a 79.38\% IQP baseline \cite{abiwardani_generalized_2025}. Parameter count does not track this pattern, however: both winning layouts use \emph{more} trainable parameters than the IQP baseline (109 and 218 against 91), not fewer. The discrete circuit ansatz applies the same principle by a different mechanism, replacing the noun's trainable box with a fixed state derived from a classical word embedding through amplitude-style state preparation, and adding only a small trainable rotation keyed to the grammatical type, specifically to handle words the parser has never seen \cite{bhatuse_resolving_2025}. Across all three, the shared premise is that a word's syntactic role, not its lexical identity, is what needs to be learned. Freezing the noun consistently improves or preserves accuracy on the tasks reported. It also cuts parameter count in SimpleSA specifically, but GeneralQC's best layouts show this need not hold in general: a redesigned ansatz can spend more parameters elsewhere in the circuit while still gaining from non-parameterized nouns.

A third cluster asks the ansatz question directly, by fixing the diagram and swapping the gate template. The widest comparison pits \textsc{lambeq}'s circuit ansätze, IQP, StronglyEntangling, Sim14, and Sim15, against its tensor-network ansätze, MPS, Spider, and plain Tensor, and finds Sim14 and Spider narrowly ahead of IQP among the circuit and tensor families respectively, with Sim15 essentially matching IQP and only StronglyEntangling clearly behind \cite{castillo_comparative_2025}. A second comparison of IQP against the Sim-et-al.\ ``circuit 14/15'' and StronglyEntangling ansätze finds the opposite depth dependence: StronglyEntangling is ahead at shallow depth, but IQP and Circuit14 converge faster and overtake it as circuit depth increases \cite{metawei_topic-aware_2023}. A related four-way comparison across six datasets gives IQP the best mean accuracy overall, while StrongEntangling wins on the hardest of them \cite{peral-garcia_comparing_2024}. The same authors elsewhere present Sim14 and the Spider ansatz only as illustrative options, without running the comparison \cite{peral-garcia_using_2024}. A narrower, non-\textsc{lambeq} comparison pits IQP against Sim14 directly on SQL circuits and finds Sim14 slightly more accurate but harder to train, concluding that IQP's simplicity is the better default for that task \cite{uotila_sql2circuits_2025}. Read together, these head-to-head studies complicate rather than confirm the workhorse status IQP holds by default elsewhere in the corpus. IQP is competitive rather than dominant: a hardware-efficient or tensor-network alternative beats it on accuracy or convergence speed in most, though not all, of the papers that run the comparison. It is also the ansatz that pulls ahead as circuit depth grows in the one study that varies depth explicitly. Its dominance in the rest of the subcorpus therefore reflects toolkit convenience at least as much as a demonstrated performance advantage.

Two studies leave the gate template unspecified. One routes the diagram through ZX-calculus rewriting before compiling it to a circuit, but delegates the compilation to a generic diagram-to-circuit call without naming the resulting gates \cite{sreedhar_towards_2025}; the only concrete claim recoverable from the paper is that ZX-simplification reduces the qubit count needed for a given sentence. The other builds a custom, non-\textsc{lambeq} pipeline and shows circuit diagrams with parameterized boxes at each atomic type, without labeling the gate template they instantiate \cite{liu_research_2024}. Neither paper permits the ansatz to be classified beyond ``some parameterized construction,'' which is itself informative: it marks the two studies in the subcorpus with the least methodological transparency about the one design choice this section is organized around.

Read together, the split shows that syntactic/DisCoCat encoding is more homogeneous at the ansatz level than at the application level: a single IQP-derived template, inherited from the family's foundational papers and packaged as \textsc{lambeq}'s default, is adopted essentially unmodified by two-thirds of the subcorpus regardless of whether the downstream task is sentiment, translation, or music. The two genuine departures pull in different directions: one strips the ansatz down by freezing nouns to cut parameters, the other benchmarks it against hardware-efficient and tensor-network alternatives and repeatedly finds it beatable, and together they show that the field's default gate template is a convenience inherited from tooling rather than a demonstrated optimum. What the ansatz choice does not touch, in any of the four buckets, is qubit count, which stays governed by the grammar diagram itself \cite{sreedhar_towards_2025,uotila_sql2circuits_2025}.

\subsection{Structured Feature-Map Encoding Across Constructions}
\label{subsec:structured_across_constructions}

\begin{table*}[t]
\caption{Structured feature-map encoding split by map construction in the final primary corpus.}
\label{tab:structured_by_construction}
\centering
\scriptsize
\begin{tabular}{@{}p{0.17\textwidth} p{0.06\textwidth} p{0.13\textwidth} p{0.24\textwidth} p{0.26\textwidth}@{}}
\hline
\textbf{Map construction} & \textbf{Studies} & \textbf{Median qubits (range)} & \textbf{Interactions encoded} & \textbf{Host families and role} \\
\hline
ZZ / IQP maps & 4 & 10 (5--20) & Pairwise, second-order feature products via entangling phase gates & Mostly kernel geometry; one hybrid encoding layer \\
General Pauli maps & 3 & 11 (6--14) & Higher-order and non-$Z$ Pauli-string interactions & Kernel and VQC; gate combination chosen by search \\
Unnamed structured maps & 2 & not reported & Unspecified feature interactions & QSVM branch; transformer feature stage \\
\hline
\end{tabular}
\end{table*}

Structured feature-map encoding is the smallest of the four families addressed here, with nine studies, and it repeats the concentration seen for syntactic encoding: six of the nine are quantum-kernel methods, with only single instances in VQC, hybrid neural, and quantum-transformer models. Because model family again barely varies, the dimension that plays its role is the feature-map construction itself. That axis is natural for this family in any case. Unlike angle and amplitude encoding, which load feature \emph{values}, a structured map builds feature \emph{interactions} into the circuit through phase rotations and entangling gates; the choice of map is therefore a modeling assumption about which interactions matter. Two properties frame the split. Like angle encoding, the family depends on aggressive classical compression before the circuit: the high-dimensional text representation is reduced by PCA or a projection layer first. Unlike angle encoding, register width then tracks the number of retained features at roughly one qubit each. This is why the reported registers are comparatively wide, with a median of $10.5$ qubits and a range of $5$ to $20$. As shown in Table~\ref{tab:structured_by_construction}, the studies are grouped by the map they use.

ZZ and IQP maps are the default construction, accounting for four of the nine studies. These are second-order maps that entangle pairs of features through data-dependent phase gates, and they are most at home in quantum kernels, where the encoded state defines the similarity geometry that the kernel measures: a ZZFeatureMap QSVM over a twenty-qubit register \cite{widdows_near-term_2024}, an IQP feature map applied to reduced AraBERT embeddings \cite{djemmal_arabert-qc_2025}, and IQP/ZZ maps feeding projected quantum kernels \cite{keith-norambuena_quantum_2026}. The same construction also appears once outside a kernel, as the primary quantum encoding layer of a hybrid model whose BERT/T5 embeddings are projected into a five-qubit ZZ map \cite{baik_neural_2025}, which illustrates the family's borderline status: even its standard map is used both as a kernel definition and as a swappable layer inside a trainable network.

General Pauli maps, in three studies, are the more expressive and more experimental construction. These use higher-order or non-$Z$ Pauli strings and treat the exact gate content as a design choice to be tuned. A question-answering VQC selects a specific $X$/$Y$/$ZZ$ combination after a small search \cite{katyayan_supervised_2023}. A quantum-enhanced SVM adopts a Pauli-$Y$/$YYY$ feature map with a rotation factor over fourteen qubits produced by PCA \cite{ruskanda_quantum-enhanced_2023}, alongside an earlier Pauli-map sentiment kernel \cite{mohammad_masum_hybrid_2023}. This is exactly where the family's central risk becomes visible: because the map fixes which feature interactions are represented, performance depends on whether that hand-chosen structure happens to match the linguistic task, and the map is a strong assumption rather than a learned one. The remaining two studies report only an unnamed ``structured feature map'', one as a QSVM branch compared against a VQC \cite{sunkavalli_experimentation_2024} and one as a feature stage inside an aviation-text transformer \cite{jagatheesan_quantum_2026}, a reporting gap that makes them hard to place against the named constructions.

Taken together, structured feature maps are the most-benchmarked yet least-dominant of the established families. Because the map is usually presented as one candidate encoding among several, the studies that use it tend to compare it directly with the vector-loading families. Those comparisons do not favor it: the group behind the Pauli-$Y$ kernel above later reports, in the controlled study cited in the angle-encoding analysis, that a simple $R_y$ angle map outperforms a Pauli feature map under matched circuits \cite{ruskanda_evaluating_2025}. The concentration in kernels is therefore the coherent core of the family, the one setting where the map is the model rather than a replaceable layer, while the scattered VQC, hybrid, and transformer instances are experiments in borrowing a designed geometry into a trainable pipeline. This leaves the family with a narrow but principled niche, testing whether an engineered similarity space helps for text, subject to a double dependency the other families do not share: its usefulness rests on both the feature-map design and the compression stage that precedes it.

\subsection{Emerging Families}
\label{subsec:emerging_specialized_interfaces}

The residual group of eight studies comprises the three emerging families of the taxonomy, and unlike the four major families it cannot be split by model family for a different reason: it is already almost mono-cultural in architecture. Seven of the eight studies are variational or recurrent circuits, with a single hybrid-neural case and no kernel, compositional, or transformer instances at all. What varies, and what organizes this subsection, is the specialized mechanism itself, because each one is best understood as a targeted patch for a specific limitation of the mainstream encodings analyzed above rather than as a general-purpose alternative to them. As shown in Table~\ref{tab:specialized_by_mechanism}, the studies are grouped by mechanism, naming the trade each one makes.

\begin{table*}[t]
\caption{Emerging encoding families split by mechanism in the final primary corpus.}
\label{tab:specialized_by_mechanism}
\centering
\scriptsize
\begin{tabular}{@{}p{0.16\textwidth} p{0.06\textwidth} p{0.12\textwidth} p{0.28\textwidth} p{0.24\textwidth}@{}}
\hline
\textbf{Mechanism} & \textbf{Studies} & \textbf{Median qubits (range)} & \textbf{Core mechanism} & \textbf{Host families and trade} \\
\hline
Data re-uploading & 4 & 7 (6--18) & Same features re-injected across circuit depth, alternating with trainable gates & VQC, recurrent; trades qubit width for depth \\
Hybrid embedding & 3 & 6 (4--8) & A classical encoder emits a compact latent vector consumed by a quantum layer & VQC, hybrid; trades quantum role for classical preprocessing \\
Hamiltonian/observable & 1 & 9 & Data mapped to Pauli-string Hamiltonians or observables, with no state preparation & VQC; trades state encoding for observable encoding \\
\hline
\end{tabular}
\end{table*}

Data re-uploading, introduced in Section~\ref{subsec:other_encoding_interfaces}, is best read as the response to the width problem that constrains angle and amplitude encoding: re-injecting the same feature stream across circuit depth, rather than loading it once, lets a small register carry more information than a single loading step would allow. Its four studies span a six-to-eighteen-qubit range, wider than angle encoding's fixed small band but nowhere near amplitude encoding's $\log_2$-scale registers, because here the burden is shifted into depth rather than width.

Hybrid embedding is the response to the high-dimensionality problem, and by the same logic it carries the smallest registers in the emerging group, a median of six qubits (Section~\ref{subsec:other_encoding_interfaces}): a classical encoder absorbs the heavy compression before a small quantum layer ever sees the result. This is the same attribution problem already flagged for hybrid-neural angle encoding (Section~\ref{subsec:angle_across_families}), taken to its limit: the quantum tail's independent contribution is the hardest of any interface to isolate.

Hamiltonian and observable encoding, the single study covered in Section~\ref{subsec:other_encoding_interfaces}, is the response to the state-preparation problem that undermines amplitude encoding's qubit savings (Section~\ref{subsec:amplitude_across_families}): it reads text-derived information out through observables rather than preparing a state at all, sidestepping state preparation entirely while keeping a logarithmic, nine-qubit budget. With only one study, it remains an idea rather than an established interface.

Taken together, the emerging families confirm the reading of the major families from the outside in: each is defined by the mainstream limitation it tries to escape, data re-uploading trading width for depth, hybrid embedding trading quantum responsibility for classical preprocessing, and Hamiltonian encoding trading state preparation for observables. Their shared confinement to variational and recurrent circuits shows that this experimentation is happening entirely within the trainable-circuit world, not in the kernel or compositional families whose encodings are structurally fixed. The evidence, however, is thinner and less comparable than for any established family: these studies differ in encoding and in where the quantum component sits and how much work precedes it. The natural questions (whether depth is worth reduced width, whether the quantum layer adds value beyond the classical encoder, and whether observable encoding beats state preparation) remain open on the current corpus.

\section{Temporal Evolution of Encoding and Model Choices}
\label{sec:temporal_evolution}

Every family and mechanism examined so far has been analyzed as a cross-section of the corpus, encoding against model family, encoding against ansatz, at a single frozen point. This section reintroduces the one dimension held constant throughout, publication year, and asks how large the field has grown and why its composition changed as it grew, drawing on the mechanisms already identified in Sections~\ref{sec:taxonomy} and~\ref{sec:comparative_analysis} to interpret the trend rather than merely report it. The final corpus spans six publication years, from two studies in 2021 to a peak of 41 in 2025. The last calendar year is included at 27 studies but was not complete at the time of writing, since the search preceded the year's close. Any apparent 2026 decline is therefore a search-cutoff artifact rather than a documented contraction, and it is treated as such throughout.

\subsection{Growth and the Changing Encoding Mix}
\label{subsec:temporal_encoding}

\Figure[t!](topskip=0pt, botskip=0pt, midskip=0pt)[width=\textwidth]{figures/encoding_by_year.pdf}{Encoding family by publication year in the final primary corpus (2026 is a partial year at the time of the search).\label{fig:encoding_by_year}}

Figure~\ref{fig:encoding_by_year} decomposes the field's fivefold growth by encoding family, and the shape of that growth points to a specific mechanism rather than to uniform diversification. Syntactic/DisCoCat encoding, the only family represented in 2021, accounted for its entire output that year and still half of it in 2022, but its share falls monotonically after that, to 38\% in 2023, 19\% in 2024, and 11\% by 2026, even while its absolute count stays roughly flat at three to eight studies a year. Angle and amplitude encoding absorb almost all of the volume that arrives after 2023. Angle rises from a quarter to roughly half of the yearly corpus. Amplitude, after an early spike to a quarter of the small 2022 corpus and a near-total drop to under 5\% in 2023, settles into a stable quarter to third of every year from 2024 onward.

The most plausible driver of this split is a difference in barrier to entry rather than in explanatory adequacy. This is the review's own interpretation of the shift, offered because the corpus-wide counts support it structurally; no paper in the corpus frames its encoding choice explicitly in these terms, so the mechanism should be read as an inference from the aggregate pattern, not a documented finding. A DisCoCat study requires committing to a grammar formalism, a parser capable of producing pregroup or CCG derivations, and the narrow \textsc{lambeq}/\textsc{DisCoPy} toolchain that Section~\ref{subsec:syntactic_across_ansatze} showed also constrains the ansatz. The payoff is a linguistically interpretable wiring, but the cost is paid again on every new sentence, language, or domain. This is why the same section found the family's practicality to fall away the further the input drifts from a short, curated English sentence. Angle and amplitude encoding impose no comparable commitment: either accepts any fixed-length numeric vector. Once general-purpose pretrained text embeddings became a routine input, available to a researcher with no linguistic-formalism background at all, the marginal cost of a new study dropped to whatever off-the-shelf VQC or kernel recipe was already circulating in the broader quantum-machine-learning literature. Section~\ref{subsec:amplitude_across_families} located that shift concretely: every VQC-hosted amplitude study from 2024 onward runs on dense transformer embeddings, and 2024 is the year amplitude's yearly share stabilizes in Figure~\ref{fig:encoding_by_year}. The remaining families fit the same account from the other side. Structured feature-map encoding stays in the single digits every year, apart from a one-off spike to 14\% in 2023, and otherwise never moves far from three to eight percent. The three emerging families (data re-uploading, hybrid embedding, and Hamiltonian/observable encoding) together never exceed one-tenth of any year's output. Each still demands bespoke, non-templated engineering (a decomposed recurrent layer, a co-designed classical encoder, a custom observable construction) that scales with the number of researchers willing to do that engineering rather than with the size of the field, so their share cannot expand the way a copy-paste encoding's can. The field's growth, in short, has followed the path of least engineering resistance: the encodings that ask the least of a newcomer, not the ones judged most linguistically or computationally motivated in Section~\ref{sec:taxonomy}.

\subsection{From Compositional Monoculture to Architectural Pluralism}
\label{subsec:temporal_model_family}

\Figure[t!](topskip=0pt, botskip=0pt, midskip=0pt)[width=\textwidth]{figures/model_family_by_year.pdf}{Model family by publication year in the final primary corpus (2026 is a partial year at the time of the search).\label{fig:model_family_by_year}}

Figure~\ref{fig:model_family_by_year} tells the same story from the architecture side, and because compositional QNLP and syntactic encoding are the same 29 studies (Section~\ref{subsec:syntactic_across_ansatze}), its decline is not a second, independent observation but the identical erosion seen from a different axis: compositional models are the entire 2021 corpus, half of 2022, and down to 11\% by 2026. What replaces that share, however, is not one architecture but four, and the shape of the replacement is again best explained by portability rather than by a documented technical advantage. VQCs, the largest beneficiary, grow from an eighth of the 2022 corpus to about a third by 2025 and just under half by 2026. A VQC is not a QNLP-specific construction; Section~\ref{sec:comparative_analysis} established it as the default host for both angle and amplitude encoding precisely because it asks nothing of the input beyond a numeric vector. Its rise is the model-family expression of the same embedding-commoditization pattern identified above. Quantum transformers and quantum kernels, by contrast, hold a steady, unspectacular seven-to-fifteen percent share from 2022 onward, with quantum kernels dipping just under that band to 7.4\% in the partial 2026 count. Both require a deliberately quantum-native design decision (a state-overlap readout for kernels, an attention-like circuit structure for transformers) that cannot simply be copied from a generic VQC recipe. Their growth therefore tracks the smaller, slower-growing population of studies willing to make that architectural commitment rather than the field's overall volume.

Hybrid neural models are the one genuine outlier: negligible through 2024, then 27\% of the 2025 corpus, the single largest jump of any family in any year, before falling back to 4\% in the partial 2026 count. Because 2026 is incomplete, the fall cannot yet be read as a reversal. The 2025 spike itself is real, and it is consistent with the same low-barrier mechanism taken one step further. A hybrid neural model, by construction, puts most of the representational burden on a pretrained classical encoder and reserves only a small quantum submodule; it is the architecture that asks the least genuinely quantum design work of any family in the corpus. It is also the family Section~\ref{subsec:emerging_specialized_interfaces} flagged as hardest to attribute to the quantum component, since the classical encoder does most of the work. Read together, the two observations support the explanation already offered for the encoding-level trend: the field's fastest-growing segments are consistently those with the lowest entry cost for a researcher bringing a pretrained classical model and a small amount of quantum machinery, not those built around a bespoke quantum-native design.

\section{Task-Level Fit Between Encoding and Model Family}
\label{sec:task_fit}

The preceding sections asked which encoding a study uses and when; this one asks whether the choice ever tracks the task itself. Is there a task family whose structure gives a principled reason to prefer one encoding over another, and does the corpus show the field acting on such a reason? The question is posed structurally, in terms of what an encoding computes, rather than through reported results. The reason for that choice is set out at the end of the section: the corpus provides no basis on which the performance figures of two studies can be compared.

\Figure[t!](topskip=0pt, botskip=0pt, midskip=0pt)[width=\textwidth]{figures/encoding_by_task_family.pdf}{Encoding family by NLP task family in the final primary corpus. A study evaluating several task families is counted once in each.\label{fig:encoding_task}}

Figure~\ref{fig:encoding_task} sets the stage: sentiment and other text classification together absorb the bulk of the corpus and every encoding family with any following at all, so before asking whether encoding tracks task it is worth asking whether these two task families can answer the question. They cannot. The encoding mix inside each is essentially the corpus-wide mix: angle encoding takes 24 of the 51 sentiment studies and 21 of the 56 other-classification studies, amplitude 12 and 9, and syntactic 8 and 15, so no family's share within either task departs from its corpus-wide share by more than about seven percentage points. Nothing here suggests that classification selects for an encoding. This is not surprising in hindsight: encoding choice inside these two task families tracks the corpus-wide habits documented in Sections~\ref{sec:comparative_analysis} and~\ref{sec:temporal_evolution}, mostly which toolchain and which pretrained embedding a group already had on hand, rather than anything about sentiment or topic classification specifically. Because dozens of unrelated research groups are all applying broadly interchangeable recipes to broadly interchangeable binary and multi-class labeling problems, classification is where the corpus has the most data and the least power to say anything about task-encoding fit.

Similarity and disambiguation is the one task family where a real, mechanistically grounded pattern emerges. Amplitude encoding accounts for six of the family's thirteen studies, roughly double its share of the corpus overall, although on thirteen studies that enrichment is suggestive rather than established. The mechanism, not the count, is what carries the argument here, and Section~\ref{subsec:amplitude_across_families} already established it: an amplitude-encoded state's inner product literally is a cosine similarity, so a SWAP test or overlap measurement performs the task's core computation rather than merely feeding a downstream classifier. The same mechanism explains why three of those six sit in the quantum-kernel model family \cite{cavar_word_2025,cavar_old_2026,cavar_semantic_2024}: a kernel method is little more than a trained readout wrapped around exactly that overlap. What the fit delivers, however, is feasibility rather than superiority, and those three studies say so themselves. Their stated contribution is that the encoded state reproduces classical embedding similarity with minimal loss at logarithmic qubit cost, which is precisely what the mechanism predicts: an amplitude state loads the classical geometry it is given, so an overlap readout can reproduce a classical cosine similarity faithfully but has no channel through which to improve on it. The structural fit is real, and this is the one place in the corpus where the encoding provably does the work the task needs, but what it buys is efficiency rather than a better answer. It is telling that the field arrived at this fit by mathematical necessity rather than by systematic search: no study set out to compare encodings on similarity tasks and reached the conclusion. It falls out of what amplitude encoding is.

Relational tasks (question answering and, to a lesser extent, similarity itself) show a weaker echo of the same logic on the compositional side, but here the corpus is too thin to test anything. Syntactic/DisCoCat encoding takes two of the seven question-answering and retrieval studies and five of the thirteen similarity and disambiguation studies, which is at least consistent with the intuition that grammatical structure should matter most when a task is explicitly about how words or entities relate; on those denominators, however, the departure from the family's corpus-wide share is not evidence of anything: one study either way moves the question-answering share by fourteen percentage points. The structural argument is in any case in tension with the practical constraint documented in Section~\ref{subsec:syntactic_across_ansatze}: the family's practicality erodes precisely as a task departs from a short, curated sentence, and relational tasks are exactly where compositional structure is best motivated and exactly where the grammar-to-circuit pipeline is under the most strain. Which of the two effects dominates is not something six studies can settle.

Sequence labeling, question answering, and generation or language modeling are each covered by only six to ten studies, too few to support a claim either way, but they fail differently. Sequence labeling is almost entirely angle-encoded, five of its six studies, so there is nothing to compare: one encoding occupies the cell rather than competing against alternatives inside it. Question answering and generation, by contrast, are spread thinly across several encodings, but each individual pairing is one or two papers from a single research thread rather than a pattern a second, independent group has reproduced. Neither failure mode supports a task-encoding claim, but for opposite reasons: one cell has no variation to explain, the others have too little replication to trust.

Only three studies fall outside the six families, and it is worth saying what they are, because the residual is small by construction rather than by neglect. Two are quantum topic models \cite{kabir_qtopic_2026,qiao_quantum_2024}: both are unsupervised, evaluated by topic coherence and diversity rather than against labels, so they fit neither the supervised families nor generation; one uses angle encoding and the other amplitude, which on two studies settles nothing. The third is Quanthoven \cite{miranda_quantum_2022}, which trains DisCoCat circuits to classify and then compose music; it is the one corpus member whose task is not linguistic at all, retained because its grammar-to-circuit pipeline is the same construction the compositional family uses on text.

This section's central finding is therefore a methodological one rather than a performance-based one: it is what licenses the structural reading used throughout. The corpus cannot be pooled on results. Each study selects its own dataset, its own classical comparison, and its own evaluation protocol: 141 distinct datasets are named across the 125 studies, and 122 of them are used by exactly one paper (Section~\ref{subsec:reproducibility_gap}); the classical systems compared against are nearly unique to each paper as well; and 15 studies report no classical baseline at all. A reported accuracy figure in this corpus is thus a statement about one pipeline, on one dataset, against one baseline; the difference between two such figures, drawn from two different studies, is not attributable to the encoding that distinguishes them. This is why the one task--encoding fit that the corpus does support has to be argued from mechanism rather than measured from results, and why even that fit resolves into efficiency rather than superiority. Determining whether an encoding does better on a task, rather than merely fitting it, requires the shared benchmarks called for in Section~\ref{subsec:future_agenda}.

\section{Experimental Rigor: Simulation, Hardware, and Reproducibility}
\label{sec:experimental_rigor}

The preceding sections asked what is encoded and by which model; this section asks how seriously the corpus tests any of it. Every primary study was coded for its execution regime (noiseless simulation, noisy simulation, real hardware, or a mixed protocol), the simulation framework and backend used, code availability, and whether scalability was discussed. These fields describe the experimental apparatus underneath the results of Sections~\ref{sec:comparative_analysis}--\ref{sec:task_fit} and support a question those sections cannot answer: not which encoding is used, but how much evidence backs the claim that it works.

\subsection{Execution Regimes Across the Corpus}
\label{subsec:execution_regimes}

Noiseless statevector simulation dominates: 96 of 125 studies (76.8\%) report exact simulation only. A further 16 (12.8\%) simulate under an explicit noise model, and 11 (8.8\%) report a mixed protocol pairing simulation with a hardware validation run. Only two studies (1.6\%) run exclusively on real hardware with no simulated baseline: a quantum-kernel KNN classifier \cite{zeguendry_quantum_2024} and a quantum self-attention model on a 72-qubit superconducting processor \cite{he_training_2024}. Combining the mixed and hardware-only categories, 13 studies (10.4\%) are \emph{hardware-touched}; widening this to any engagement with noise (noisy simulation, mixed protocol, or hardware-only) raises the count to 29 studies (23.2\%), termed \emph{noise-or-hardware engagement}. The remaining three-quarters of the corpus never leaves a noiseless simulator. Hardware-touched status required the study to name the device it ran on: three studies that report a quantum backend but resolve, on inspection, to a fake-hardware noise model or a vendor simulator are coded here as noisy or noiseless simulation rather than as hardware runs.

\begin{table}[t]
\caption{Execution regime by encoding family in the final primary corpus.}
\label{tab:execution_by_encoding}
\centering
\scriptsize
\begin{tabular}{@{}lcccc@{}}
\hline
\textbf{Encoding family} & \textbf{Noiseless} & \textbf{Noisy sim} & \textbf{Mixed} & \textbf{HW-only}\\
\hline
Angle & 42 & 6 & 2 & 1\\
Amplitude & 24 & 2 & 1 & 1\\
Syntactic/DisCoCat & 19 & 6 & 4 & 0\\
Structured feature map & 5 & 1 & 3 & 0\\
Data re-uploading & 2 & 1 & 1 & 0\\
Hybrid embedding & 3 & 0 & 0 & 0\\
Hamiltonian/observable & 1 & 0 & 0 & 0\\
\hline
\end{tabular}
\end{table}

Table~\ref{tab:execution_by_encoding} crosses execution regime with the taxonomy of Section~\ref{sec:taxonomy}. As a whole the table shows no clean split between the two dimensions, and four of the seven families contribute fewer than ten studies each, so the per-family percentages below are description rather than demonstrated difference, except for the one contrast at the end of this paragraph that rests on denominators large enough to carry weight. Descriptively, engagement with noise and hardware concentrates in the families Section~\ref{sec:comparative_analysis} already identified as small and toolchain-bound. Syntactic/DisCoCat encoding is hardware-touched somewhat more often than the corpus average (13.8\% against 10.4\%) and engages noise or hardware at more than twice the rate of angle and amplitude encoding (34.5\% against 16.5\%), because its \textsc{lambeq}/\textsc{DisCoPy}/\textsc{pytket} stack (Section~\ref{subsec:syntactic_across_ansatze}) ships with direct IBM Quantum integration and its shallow, few-qubit circuits are cheap to queue on real devices \cite{lorenz_qnlp_2023,meichanetzidis_grammar-aware_2023,peral-garcia_comparing_2024,miranda_quantum_2022}. Structured feature maps and data re-uploading show the highest hardware- and noise-touched shares of all, but on denominators of nine and four studies, so these describe a handful of papers rather than a trend. The contrast that does hold: angle and amplitude encoding, 79 of the 125 studies, engage noise or hardware at less than half the rate of the rest of the corpus (13 of 79, 16.5\%, against 16 of 46, 34.8\%). Restricted to real-device runs the asymmetry points the same way, on smaller numbers (5 of 79, 6.3\%, against 8 of 46, 17.4\%). Since these two encodings carry the field's growth (Section~\ref{subsec:temporal_encoding}), most of their reported accuracy numbers (66 of 79) are simulator artifacts, numbers silent on whether the same circuit would survive decoherence, gate error, or readout noise on real hardware.

Hybrid embedding and Hamiltonian/observable encoding report no hardware or noisy-simulation validation whatsoever: every study in both families is exact-simulation only. For Hamiltonian encoding this is one study, so no pattern can be claimed, but it means the family's central promise, that observable-based readout sidesteps amplitude encoding's state-preparation cost (Section~\ref{subsec:emerging_specialized_interfaces}), has not been tested under any condition that would actually stress it.

\subsection{Execution by Model Family and the Recurrent-Network Blind Spot}
\label{subsec:execution_by_model_family}

The same asymmetry appears when execution is crossed with model family instead of encoding. Quantum kernels are the most hardware-touched family (3 of 13 run on a named device, 23.1\%), consistent with their concentration in amplitude and structured-feature-map encoding (Section~\ref{subsec:amplitude_across_families}); on the wider measure compositional models lead instead, 34.5\% of them engaging noise or hardware against the kernels' 30.8\%, the same figure as syntactic encoding, since the two categories are the same 29 studies. VQCs, hybrid-neural models, and quantum transformers cluster around 19--27\%. Recurrent quantum models are the one family with \emph{zero}: all 12 recurrent-network studies (QLSTM, QRNN, and related sequential architectures) report noiseless simulation only. This matters because a recurrent quantum cell re-applies its encoding and ansatz once per token (Section~\ref{subsec:angle_across_families}), so circuit depth, and exposure to decoherence and gate error, scales with sequence length in a way a single-shot VQC does not: the family whose central architectural claim rests most on reusing a small register cheaply is the one for which noise robustness has been tested least.

\subsection{Hardware Validation Peaked in 2024 and Then Collapsed}
\label{subsec:hardware_timeline}

Section~\ref{subsec:temporal_encoding} showed the corpus grew roughly twenty-fold between 2021 and 2025, but the hardware-touched share does not grow with it. It climbs slowly to a peak and then collapses: 0\% in 2021, 12.5\% in 2022, 14.3\% in 2023 and 26.9\% in 2024, then 2.4\% in 2025 and 3.7\% in 2026 (partial year). Noisy-simulation share does not follow the collapse, holding between 12\% and 15\% in every year from 2022 onward except 2024, so the late-period loss is specific to real-device runs. Pooling the window puts the size of the drop plainly: 11 of 57 studies published 2021--2024 are hardware-touched (19.3\%), against 2 of 68 in 2025--2026 (2.9\%). The matching contrast for noise-or-hardware engagement generally is far milder, 17 of 57 (29.8\%) against 12 of 68 (17.6\%), consistent with real-device execution specifically falling away rather than noise modeling. Setting aside the two-paper 2021 cohort, the two most recent years thus show the lowest combined rate of hardware or noise engagement of the entire window, 17.1\% and 18.5\% against 25--35\% across 2022--2024, at the exact moment the field's output peaked, the experimental-rigor counterpart to the low-barrier mechanism Section~\ref{subsec:temporal_encoding} proposed for the encoding mix: growth ran through the protocol that asks the least of a newcomer, a local noiseless simulator call rather than a queued hardware job or a calibrated noise model. Two studies reach hardware anyway: a neural-quantum-embedding pipeline validated on an IBM backend \cite{baik_neural_2025}, and a quantum-secure federated-learning system validated on a 156-qubit backend, though motivated by communication security rather than an NLP-performance claim \cite{liu_experimentally_2026}.

\subsection{Framework Fragmentation and the Reproducibility Gap}
\label{subsec:reproducibility_gap}

Even where a study stays in simulation, the tooling underneath it is far from standardized. 37 of 125 studies (29.6\%) do not name the simulation framework at all, the single largest category. Among those that do, PennyLane is most common (21 sole and 22 more in combination, 43 studies in all), Qiskit is second (9 sole, 25 mentions in all), and \textsc{lambeq} anchors the syntactic/DisCoCat subcorpus (Section~\ref{subsec:syntactic_across_ansatze}); the rest (TensorCircuit \cite{liu_hcqnsc_2026}, Amazon Braket Hybrid Jobs, Qulacs, TensorFlow Quantum, Flax) are each unique to a single paper. Concentration in the corpus is therefore at the framework layer alone: outside the \textsc{lambeq} cluster and a PennyLane/Qiskit core, the stacks built on top of those frameworks are close to unique per paper. This compounds a second gap: only 25 of 125 studies (20.0\%) report publicly available code. Between the two gaps, 34 studies (27.2\%) report neither code nor a framework and only 22 (17.6\%) report both, so a substantial fraction of the corpus cannot be independently re-run even in simulation, let alone re-validated on hardware.

Scalability discussion fares better on paper: 91 studies (72.8\%) discuss it against 34 that do not. But like the evaluation protocols of Section~\ref{sec:task_fit}, this is self-reported against whatever standard each paper's authors chose, ranging from a formal resource analysis to a single sentence deferring qubit count to future work. Combined with the noiseless-simulation dominance of Section~\ref{subsec:execution_regimes}, it remains unclear whether any encoding in Section~\ref{sec:taxonomy} would still work outside a simulator at a larger register size, generalizing the state-preparation caution Section~\ref{subsec:amplitude_encoding} raised specifically for amplitude encoding.

\section{Trends and Open Challenges}
\label{sec:trends_challenges}

The preceding five sections examined the corpus one angle at a time: what is encoded (Section~\ref{sec:taxonomy}), how the same encoding is repurposed across model families (Section~\ref{sec:comparative_analysis}), how the field's composition changed over time (Section~\ref{sec:temporal_evolution}), whether the choice ever tracks the task (Section~\ref{sec:task_fit}), and how much evidence backs any of it (Section~\ref{sec:experimental_rigor}). This section reads them together, states plainly what the encoding bottleneck named in the abstract is, and converts the gaps into an agenda for future work.

\subsection{Where the Field Is Heading}

Three trends established separately above describe a single movement. The field is consolidating on the encodings that ask least of a newcomer, angle and amplitude, at the expense of the grammar-dependent syntactic/DisCoCat tradition (Sections~\ref{subsec:temporal_encoding}--\ref{subsec:temporal_model_family}). That consolidation is landing on simulators rather than on hardware: 2025 and 2026, the years the corpus grew fastest, carry the lowest combined rate of hardware or noisy-simulation engagement of the entire six-year window (Section~\ref{subsec:hardware_timeline}), and the shortfall concentrates in the same two families that carry the growth (Section~\ref{subsec:execution_regimes}). And the corpus's one mechanistically grounded task--encoding fit, amplitude encoding on similarity and disambiguation, is small, drawn disproportionately toward quantum kernels, and arrived at by mathematical necessity rather than by systematic search (Section~\ref{sec:task_fit}).

These are not three independent observations: the field's growth curve and its evidentiary curve are decoupling. Publication volume is rising fastest in exactly the segment least equipped to demonstrate anything about noise, hardware behavior, or a genuine quantum contribution, while the segment with the clearest structural argument for one remains a slow-growing niche. Underneath sits the incommensurability documented in Section~\ref{sec:task_fit}, which leaves a reported figure carrying little evidential weight beyond the single pipeline that produced it. None of these gaps belongs to one encoding or one model family; they are properties of how the field as a whole reports its experiments, which is why closing them is a precondition for the task--encoding questions to be answerable at all.

\subsection{The Encoding Bottleneck}
\label{subsec:encoding_bottleneck}

Every family reviewed in Sections~\ref{sec:taxonomy} and~\ref{sec:comparative_analysis} but one is, underneath its individual mechanism, a different response to the same arithmetic. Let $d$ denote the dimensionality of the text representation a study starts from, in this corpus from a handful of TF-IDF or handcrafted features up to the 768 or more of a transformer embedding, and let $q$ denote the qubit budget it actually uses, empirically two to sixteen across the angle and amplitude subcorpora (Tables~\ref{tab:angle_by_family} and~\ref{tab:amplitude_by_family}). Angle encoding spends one rotation on each retained feature, so the register must be about as wide as the feature count $d' \leq d$ that survives a classical reduction applied \emph{before} the circuit is built: the mismatch is resolved by discarding information, and the denser variants of Section~\ref{subsec:angle_encoding} narrow the register only by paying in loading depth. Amplitude encoding instead packs the full vector into $q = \lceil \log_2 d \rceil$ qubits, but preparing an arbitrary $2^{q}$-dimensional amplitude vector generically requires a gate count scaling with $d$ rather than with $q$ \cite{mottonen2005transformation,plesch2011quantum} (Section~\ref{subsec:amplitude_encoding}); here the mismatch is paid in circuit depth, a cost Section~\ref{subsec:amplitude_across_families} showed is routinely absorbed into a simulator call and left out of the reported qubit budget. Structured feature maps and data re-uploading sit between the two, trading some width for some depth without eliminating either cost, and hybrid embedding relocates the mismatch upstream into a classical encoder rather than resolving it. Syntactic/DisCoCat encoding alone sidesteps the arithmetic instead of trading inside it: $q$ is set by the grammatical structure of the sentence, not by the dimensionality of an embedding (Section~\ref{subsec:syntactic_encoding}), which is why Section~\ref{sec:taxonomy} treats it as categorically different and why it is the one encoding whose qubit count Section~\ref{subsec:syntactic_across_ansatze} found untouched by ansatz choice.

This is the encoding bottleneck named in the abstract: a persistent gap between $d$, set by what classical NLP representations require to be informative, and $q$, set by what near-term hardware can support. Every non-syntactic family is best understood as a choice of \emph{which cost absorbs the gap}, lost information, added depth, or a relocated classical burden, rather than as a choice that closes it, and no study in the corpus reports an encoding that escapes all three. The bottleneck is structural, not a matter of engineering effort within any single family.

\subsection{Open Challenges}
\label{subsec:open_challenges}

The gaps documented throughout the review reduce to five open challenges, each derived from the corpus itself rather than imported from general quantum-computing commentary.

\begin{enumerate}
\item[C1)] \emph{The encoding bottleneck and the information it costs.} Outside the syntactic family the reduction stage discards most of the representation before the circuit ever sees it, 768-dimensional embeddings cut to ten \cite{ahmad_hybrid_2023}, four \cite{kumar_hybrid_2026}, or two \cite{altintas_beyond_2025} features, while state-preparation costs are routinely left out of the reported qubit budget. The corpus's mitigations, divide-and-conquer preparation \cite{zhang_entangled_2024}, observable encoding \cite{tiblias_efficient_2026}, and depth-for-width trades \cite{silver_lexiql_2024}, change which cost is paid without closing the gap, and no surveyed study quantifies what the reduction removes.

\item[C2)] \emph{Incommensurable evaluation.} Each study chooses its own dataset, classical comparison, and evaluation protocol; 15 report no classical baseline at all, and among those that do the compared system is very nearly unique to each paper. Even the one structurally motivated pairing, amplitude encoding on similarity tasks, is presented by its own studies as reproduction of classical similarity rather than an improvement on it \cite{cavar_word_2025,cavar_old_2026}. Open because no shared benchmark with fixed classical baselines exists against which results from different studies could be made commensurable.

\item[C3)] \emph{The simulation--hardware gap.} 96 of 125 studies never leave a noiseless simulator; real-device validation peaked in 2024 and then collapsed to at most 3.7\% in 2025--2026 (Section~\ref{subsec:hardware_timeline}), the only hardware-validated studies of those two years being \cite{baik_neural_2025,liu_experimentally_2026}, the latter reaching hardware for reasons orthogonal to an NLP-performance claim. Open because noiseless simulation is silent on decoherence, gate error, and readout noise, precisely the properties that decide near-term viability.

\item[C4)] \emph{Absent scaling evidence.} 102 of 125 studies evaluate a single register width; only 23 test more than one, and systematic sweeps such as the 1--16-qubit comparisons of \cite{zhang_entangled_2024,kim_quantum_2025} are rare. Open because every resource argument in Section~\ref{subsec:encoding_bottleneck} is asymptotic while the evidence is single-point.

\item[C5)] \emph{Reproducibility and benchmark triviality.} 100 studies (80.0\%) release no code and 37 (29.6\%) do not name their simulation framework (Section~\ref{subsec:reproducibility_gap}); of 141 distinct datasets, 122 are used by exactly one paper and only 7 recur in four or more. Those that do recur are either generic sentiment benchmarks carrying no QNLP-specific signal (IMDb, 12 papers; Amazon reviews, 8; Yelp, 7) or the $\sim$100-sentence synthetic MC/RP toy corpora built into \textsc{lambeq} (12 and 9 papers), on which near-perfect accuracies are reported \cite{lorenz_qnlp_2023,bouakba_ensemble_2023,liu_research_2024}. Open because at this reporting standard neither independent re-execution nor a shared benchmark is possible.
\end{enumerate}

\subsection{An Agenda for Future Work}
\label{subsec:future_agenda}

Four actions follow directly from C1--C5.

\emph{Noise-in-the-loop evaluation as a default.} Report accuracy under an explicit, calibrated noise model as a minimum bar for publication, rather than treating noise as a qualitative caveat. The case follows from C3: angle and amplitude encoding are simultaneously the fastest-growing families and the least hardware- or noise-validated, and recurrent quantum architectures, which re-apply their encoding once per token, have no noisy-simulation or hardware study anywhere in the corpus (Section~\ref{subsec:execution_by_model_family}).

\emph{Controlled, cross-family encoding comparisons.} Hold model family, dataset, and circuit depth fixed while swapping only the encoding or ansatz, making the controlled comparison the template rather than the exception. Section~\ref{sec:comparative_analysis} showed encoding and model family to be so tightly coupled that comparing ``encoding families'' is often really comparing the architectures they are bound to; the few studies that already run such comparisons are disproportionately informative relative to their count.

\emph{Shared benchmarks with fixed classical baselines.} Establish a small set of shared, encoding-agnostic NLP benchmarks with fixed classical baselines, of the kind proposed piecemeal but never standardized in the surveyed studies. This addresses C2 directly, and would let a future meta-analysis attribute a performance gap to the encoding rather than to the baseline an individual group selected.

\emph{From structural fit to a measurable gain.} Search deliberately for task families where an encoding's mathematical structure performs the task's core computation, and prefer fits where that computation is classically expensive: whether the compositional structure of syntactic encoding can be matched to relational tasks such as question answering, for instance, where Section~\ref{sec:task_fit} found a theoretical fit the corpus is too thin to evaluate. The one documented fit yields efficiency rather than a better result, since the encoded state reproduces exactly the classical geometry it loads. A targeted study, rather than another general-purpose classifier, could establish whether a structural fit exists anywhere that also produces a gain measurable against a fixed classical baseline.

\section{Conclusion}
\label{sec:conclusion}

This article treated text encoding as the primary variable of QNLP research. Under a protocol fixed before the search was executed (Section~\ref{sec:method}), it screened 376 database records and 1222 snowball candidates down to a final corpus of 125 primary studies published between 2021 and 2026. It organized them under a seven-family taxonomy, four major families (angle, amplitude, syntactic/DisCoCat, and structured feature-map) and three emerging families (data re-uploading, hybrid embedding, and Hamiltonian/observable), defined not by model architecture or ansatz but by the point at which text enters the quantum computation (Section~\ref{sec:taxonomy}).

Four findings recur across the analysis. First, encoding and model family are tightly coupled rather than independently chosen. Syntactic encoding and compositional QNLP coincide by construction and so are set aside, but the coupling holds across the remaining 96 studies as well: structured feature maps concentrate in quantum kernels, and only angle and amplitude encoding travel freely across architectures. Comparing ``encodings'' in the literature is therefore frequently a comparison of the architectures they are bound to (Section~\ref{sec:comparative_analysis}). Second, the field's rapid growth since 2023 has been a consolidation onto the encodings and model families that ask the least of a newcomer (angle- and amplitude-encoded VQCs and hybrid-neural pipelines built on pretrained embeddings), at the expense of the grammar-dependent compositional tradition that defined the field's first two years (Section~\ref{sec:temporal_evolution}). Third, that growth has not been matched by growth in experimental rigor: three-quarters of the corpus never leaves a noiseless simulator, real-device validation peaked in 2024 and then collapsed as the corpus grew fastest, and code availability and simulation framework are each under-reported by a wide margin (Section~\ref{sec:experimental_rigor}). Fourth, the corpus cannot be pooled on results: each study evaluates on a dataset, against a classical comparison, and under a protocol of its own choosing, and 15 studies report no classical baseline at all, so figures reported by two different studies are rarely commensurable. The only task--encoding fit the corpus supports therefore has to be argued from mechanism, and even that one, amplitude encoding on similarity and disambiguation tasks, where the encoded state's inner product performs the task's core computation directly, is presented by its own studies as faithful reproduction of classical similarity at logarithmic qubit cost rather than as an improvement on it (Section~\ref{sec:task_fit}).

Underlying all four findings is the encoding bottleneck stated in Section~\ref{subsec:encoding_bottleneck}: the persistent gap between the dimensionality of a useful text representation and the qubit budget near-term hardware can support. Every surveyed family except syntactic encoding resolves this gap by paying in one of three currencies (discarded information, added circuit depth, or a relocated classical burden), and none eliminates it. This framing is the taxonomy's main payoff. It turns a fragmented set of implementation choices into a small number of comparable trade-offs, and it grounds the agenda of Section~\ref{sec:trends_challenges}.

Three caveats condition these results. The search covered Scopus and IEEE Xplore and, although two rounds of backward and forward snowball sampling were used to recover work those databases miss, seeding from the field's surveys (Section~\ref{sec:method}) leaves primary work with no citation link to any survey outside its reach. The 2026 counts are partial, since the search was conducted before the year closed, so the final point of every year-over-year trend should be read with that in mind. And a few extracted fields, scalability discussion among them, are self-reported by the primary studies rather than independently re-derived. None of this changes the central conclusion. Text encoding is not a neutral preprocessing step in QNLP; it is the decision that fixes how much of a language sample survives contact with a near-term quantum circuit, and it deserves the same scrutiny as the ansatz, the optimizer, or the classical baseline it is compared against. The taxonomy and the classification scheme introduced here are offered as the basis for that scrutiny in future work, including the encoding-controlled comparisons, noise-in-the-loop evaluations, and shared benchmarks called for in Section~\ref{sec:trends_challenges}.

\bibliography{references}
\bibliographystyle{IEEEtran}

\EOD

\end{document}
